% =============================================================
% Pure O-Sequences in Codimension Three
% A Dependency-Ordered Formal Exposition
% =============================================================


\documentclass[11pt,reqno]{amsart}
%%%%%%%%%%%%%%%%%%%%%%%%%%%%%%%%%%%%%%%%%%%%%%%%%%%%%%%%%%%%%%%%%%%%%%%%%%%%%%%%%%%%%%%%%%%%%%%%%%%%%%%%%%%%%%%%%%%%%%%%%%%%%%%%%%%%%%%%%%%%%%%%%%%%%%%%%%%%%%%%%%%%%%%%%%%%%%%%%%%%%%%%%%%%%%%%%%%%%%%%%%%%%%%%%%%%%%%%%%%%%%%%%%%%%%%%%%%%%%%%%%%%%%%%%%%%
\usepackage{amsfonts}
\usepackage[T1]{fontenc}
\usepackage[utf8]{inputenc}
\usepackage{textcomp}
\usepackage{titlesec,amsmath,amssymb,amsthm,ragged2e,float}
\usepackage{array}
\usepackage{graphicx,setspace}
\usepackage[margin=0.5in]{geometry}
\usepackage[colorlinks=true,linkcolor=blue,citecolor=blue,urlcolor=blue]{hyperref}
\usepackage{xcolor}
\usepackage{fancyvrb}

\setcounter{MaxMatrixCols}{10}
%TCIDATA{OutputFilter=Latex.dll}
%TCIDATA{Version=5.50.0.2960}
%TCIDATA{<META NAME="SaveForMode" CONTENT="1">}
%TCIDATA{BibliographyScheme=Manual}
%TCIDATA{LastRevised=Monday, October 19, 2026 22:01:17}
%TCIDATA{<META NAME="GraphicsSave" CONTENT="32">}

\IfFileExists{tcilatex.tex}{\input{tcilatex}}{}
\DeclareUnicodeCharacter{03B9}{\ensuremath{\iota}}
\DeclareUnicodeCharacter{2200}{\ensuremath{\forall}}
\DeclareUnicodeCharacter{2203}{\ensuremath{\exists}}
\DeclareUnicodeCharacter{2218}{\ensuremath{\circ}}
\DeclareUnicodeCharacter{2264}{\ensuremath{\leq}}
\DeclareUnicodeCharacter{2265}{\ensuremath{\geq}}
\DeclareUnicodeCharacter{2192}{\ensuremath{\to}}
\DeclareUnicodeCharacter{21A6}{\ensuremath{\mapsto}}
\DeclareUnicodeCharacter{2260}{\ensuremath{\neq}}
\DeclareUnicodeCharacter{00D7}{\ensuremath{\times}}
\DeclareUnicodeCharacter{2115}{\ensuremath{\mathbb{N}}}
\DeclareUnicodeCharacter{211D}{\ensuremath{\mathbb{R}}}
\DeclareUnicodeCharacter{27E8}{\ensuremath{\langle}}
\DeclareUnicodeCharacter{27E9}{\ensuremath{\rangle}}
\providecommand{\R}{\mathbb{R}}
\providecommand{\N}{\mathbb{N}}
\providecommand{\Z}{\mathbb{Z}}
\allowdisplaybreaks
\hbadness=10000
\vbadness=10000
\vfuzz=5pt
\hfuzz=30pt
\fvset{frame=lines,framerule=0.7pt,rulecolor=\color{blue!60!black},
	framesep=5pt,fontsize=\scriptsize,xleftmargin=6pt,xrightmargin=6pt}
\RecustomVerbatimEnvironment{verbatim}{Verbatim}{}
\DefineVerbatimEnvironment{verbatimbox}{Verbatim}{
	frame=single,framerule=0.7pt,rulecolor=\color{blue!60!black},
	framesep=5pt,fontsize=\scriptsize,xleftmargin=6pt,xrightmargin=6pt}
\subjclass{\textbf{2020 Mathematics Subject Classification.}\,
	Primary 05E40, 13D40, 05A20;
	Secondary 06A07, 05A17, 13F55.}
\makeatletter
\def\@setsubjclass{\ifx\@subjclass\@empty\else\@subjclass\fi}
\makeatother
\keywords{Hilbert functions, pure $O$-sequences, log-concavity, unimodality,
	monomial ideals, order ideals, codimension three, socle degree, type,
	Macaulay's theorem}
\renewcommand{\abstractname}{\fontsize{16}{18}\selectfont Abstract}
\newcommand{\abstractfont}{\normalfont\fontsize{11}{13.2}\selectfont}
\newcommand{\abstractheadingsep}{\vspace{1em}}
\makeatletter
\renewenvironment{abstract}{\ifx\maketitle\relax
	\ClassWarning{\@classname}{Abstract should precede
		\protect\maketitle\space in AMS document classes; reported}\fi
	\global\setbox\abstractbox=\vtop \bgroup
	\abstractfont
	{\centering\abstractname\par}\abstractheadingsep
	\list{}{\labelwidth\z@\leftmargin3pc \rightmargin\leftmargin
		\listparindent\normalparindent \itemindent\z@
		\parsep\z@ \@plus\p@\let\fullwidthdisplay\relax}\item\relax
}{  \endlist\egroup
	\ifx\@setabstract\relax \@setabstracta \fi}
\makeatother
\setcounter{secnumdepth}{3}
\newcommand{\SecAbove}{2.6ex plus 0.6ex minus 0.15ex}
\newcommand{\SecBelow}{2.0ex plus 0.6ex minus 0.15ex}
\newcommand{\SubsecAbove}{2.0ex plus 0.5ex minus 0.15ex}
\newcommand{\SubsecBelow}{1.2ex plus 0.4ex minus 0.15ex}
\newcommand{\SubsubsecAbove}{2.0ex plus 0.5ex minus 0.15ex}
\newcommand{\SubsubsecBelow}{1.2ex plus 0.4ex minus 0.15ex}
\makeatletter
\titleformat{\section}{\normalfont\fontsize{16}{18}\selectfont\bfseries}
{\thesection.}{0.5em}{}
\titleformat{\subsection}{\normalfont\fontsize{14}{16}\selectfont\bfseries}
{\thesubsection}{0.5em}{}
\titleformat{\subsubsection}{\normalfont\fontsize{12}{14}\selectfont\bfseries}
{\thesubsubsection}{0.5em}{}
\titlespacing*{\section}{\z@}{\SecAbove}{\SecBelow}
\titlespacing*{\subsection}{\z@}{\SubsecAbove}{\SubsecBelow}
\titlespacing*{\subsubsection}{\z@}{\SubsubsecAbove}{\SubsubsecBelow}
\makeatother
\setlength{\arraycolsep}{3pt}
\renewcommand{\arraystretch}{1.3}
\newcolumntype{r}{>{\centering\arraybackslash\scriptsize$}m{3cm}<{$}}
\newcolumntype{l}{>{\raggedright\arraybackslash\scriptsize$}m{12cm}<{$}}
\newcolumntype{c}{>{\centering\arraybackslash\scriptsize$}m{2cm}<{$}}
\newcounter{axctr}\newcounter{posctr}\newcounter{defctr}
\newcounter{lemctr}\newcounter{thmctr}\newcounter{propctr}\newcounter{corctr}
\renewcommand{\theaxctr}{\arabic{axctr}}
\renewcommand{\theposctr}{\arabic{posctr}}
\renewcommand{\thedefctr}{\arabic{defctr}}
\renewcommand{\thelemctr}{\arabic{lemctr}}
\renewcommand{\thethmctr}{\arabic{thmctr}}
\renewcommand{\thepropctr}{\arabic{propctr}}
\renewcommand{\thecorctr}{\arabic{corctr}}
\newcommand{\swpax}[2]{\refstepcounter{axctr}\hypertarget{#1}{\mbox{\textbf{Ax\,\theaxctr}}}\label{#2}}
\newcommand{\swppos}[2]{\refstepcounter{posctr}\hypertarget{#1}{\mbox{\textbf{Pos\,\theposctr}}}\label{#2}}
\newcommand{\swpdef}[2]{\refstepcounter{defctr}\hypertarget{#1}{\mbox{\textbf{Def\,\thedefctr}}}\label{#2}}
\newcommand{\swplem}[2]{\refstepcounter{lemctr}\hypertarget{#1}{\mbox{\textbf{Lem\,\thelemctr}}}\label{#2}}
\newcommand{\swpthm}[2]{\refstepcounter{thmctr}\hypertarget{#1}{\mbox{\textbf{Thm\,\thethmctr}}}\label{#2}}
\newcommand{\tagged}[2]{\hypertarget{#1}{\mbox{\textbf{#2}}}}
\newcommand{\refto}[2]{\hyperlink{#1}{#2}}
\newcommand{\mcite}[1]{\text{\cite{#1}}}
\makeatletter
\let\orig@maketitle\maketitle
\renewcommand{\maketitle}{\begingroup
	\let\MakeUppercase\@firstofone
	\let\uppercase\@firstofone
	\orig@maketitle
	\endgroup}
\makeatother
\makeatletter
\newlength{\titleToAuthorSkip}
\newlength{\authorToAbstractSkip}
\setlength{\titleToAuthorSkip}{10pt plus 4pt minus 2pt}
\setlength{\authorToAbstractSkip}{20pt plus 6pt minus 3pt}
\renewcommand{\@maketitle}{\normalfont\normalsize
	\let\@makefnmark\relax
	\let\@thefnmark\relax
	\ifx\@empty\@subjclass\else
	\@footnotetext{\@setsubjclass}\fi
	\ifx\@empty\@keywords\else
	\@footnotetext{\@setkeywords}\fi
	\ifx\@empty\@thanks\else
	\@footnotetext{\@setthanks}\fi
	\@settitle
	\ifx\@empty\@authors\else
	\vspace{\titleToAuthorSkip}\@setauthors\fi
	\ifx\@empty\@date\else
	\vspace{6pt}\@setdate\fi
	\ifx\@empty\@abstract\else
	\vspace{\authorToAbstractSkip}\fi
	\@setabstract}
\makeatother
\pagestyle{plain}
\renewcommand{\seriesdefault}{\bfdefault}
\boldmath

\begin{document}
\title{ \bfseries{\Large \emph{Pure $O$-Sequences in Codimension Three.}}\\
[5pt] \emph{A Dependency-Ordered Formal Exposition} }
\author{ 
\bfseries{\large \hspace{1.5in} Brahimi Mahdi Tahar\newline
			\text{ \  } \newline
			\text{ \  } \hspace{0.5in}\text{ \  } University of Mohamed Boudiaf of M'sila, Algeria\newline
			\text{ \  } \newline
			\href{mailto:mahditahar.brahimi@univ-msila.dz}{mahditahar.brahimi@univ-msila.dz}}
}
\maketitle

\begin{abstract}
\bfseries
\setstretch{1.3}{\small This document is a step-indexed, dependency-ordered
presentation of eight problems on \emph{pure $O$-sequences}, following
\cite{stanley1978, bmmnz2012, zanello2024}. Every result appears only after
all results it depends on. Entries are split into postulates (unproved
classical statements used inside proofs) and axioms of the ambient
set-theoretic foundation (\textsf{ZF} with dependent choice where noted).
The document covers the prerequisite layer (monomial combinatorics, order
ideals, principal $O$-sequences, log-concavity and unimodality, Macaulay's
two-variable theorem) and the eight-problem layer built on it: an explicit
codimension-$3$ example; the type-$1$ log-concavity theorem; first-half
log-concavity in codimension $3$; unimodality and log-concavity of type-$2$
sequences in codimension $3$ --- the main result, which (subject to the
formal verification reported in §\ref{sec:verification}) resolves the
principal remaining open case identified by Zanello; eventual
non-log-concavity for every codimension $r\ge3$ and every type $t\ge5$; and
the existence of the least such threshold. The codimension-$r\ge5$ type-$2$
unimodality question (P6) is stated but not resolved. Each entry is
annotated with the choice strength of its proof: \textsf{ZF},
\textsf{ZF+DC}, or \textsf{ZFC}.}
\end{abstract}

\noindent \text{ \ } \newline

\begin{center}
\begin{minipage}[c]{0.50\linewidth}
		\vspace{0pt}
		\begin{Verbatim}[frame=single,framerule=0.7pt,
		rulecolor=\color{blue!60!black},framesep=8pt,
		fontsize=\fontsize{8}{9}\selectfont,
		formatcom=\color{red}\bfseries,
		xleftmargin=6pt,xrightmargin=6pt]
		Root theorems (nothing depends on them):
		P1   P2   P3   P4   P5   P7   P8
		(P6 is stated but open)

		Postulates (external inputs):
		pos:conv-LC
		pos:Macaulay2
		\end{Verbatim}
	\end{minipage}\hfill 
\begin{minipage}[c]{0.43\linewidth}
		\vspace{0pt}
		\begin{Verbatim}[frame=single,framerule=0.7pt,
		rulecolor=\color{blue!60!black},framesep=8pt,
		fontsize=\fontsize{8}{9}\selectfont,
		formatcom=\color{red}\bfseries,
		xleftmargin=6pt,xrightmargin=6pt]
		Deepest chains (by depth):
		L0 | P5
		L1 |   delta2-dich
		L2 |     Lem:delta2-formula
		L1 |   LU-dich
		L2 |     delta2-dich
		L1 |   Lem:LC-core-alg

		L0 | P2
		L1 |   principal-OSeq
		L2 |     pos:conv-LC
		L1 |   principal-pos

		L0 | P7
		L1 |   exists-non-lc-r3
		L1 |   embed-preserves
		\end{Verbatim}
	\end{minipage}

\medskip \noindent \text{ \ } \\[0pt]
{\small \refstepcounter{figure}%
\textcolor{red}{\textbf{Figure \thefigure.}\enspace
			Summary of root theorems, postulates, and the deepest dependency
			chains.}} \label{fig:summary}
\end{center}

\newpage \pagestyle{plain} \pagenumbering{arabic} \setstretch{1.3}

% ============================================================

\section{Introduction}

% ============================================================
\noindent The document is organised into six parts: \emph{Axioms}, \emph{%
Postulates}, \emph{Definitions}, \emph{Statements} (Lemmas, Theorems), \emph{%
Proofs}, and \emph{Scripts}. Following the reverse-mathematics convention of 
\cite{simpson2009}, entries are split into two primitive layers: \emph{%
axioms of the ambient foundation} and \emph{postulates} (unproved classical
statements admitted inside proofs). Every entry in a later part may depend
only on earlier entries, and no entry may depend on anything outside the
union of axioms and postulates. Lemmas and Theorems form a single dependency
layer: an entry may depend on an earlier Lemma or Theorem, so long as the
global order is respected.

\medskip \noindent The eight problems concern pure $O$-sequences $\Gamma
\subset \N^r$, i.e.\ finite order ideals of monomials whose maximal elements
all have a common total degree (the \emph{socle degree} $e$); the \emph{type}
is the number of maximal elements. The pure $O$-sequence is $h_\Gamma(d) =
\#\{m\in\Gamma : |m|=d\}$. Interest in these sequences originates with
Stanley's work on Hilbert functions of monomial level algebras \cite%
{stanley1978}, was formalised in the Boij--Migliore--
Mir\'o-Roig--Nagel--Zanello memoir \cite{bmmnz2012}, and the remaining open
cases were reviewed by Zanello \cite{zanello2024}. Problem P5 ---
log-concavity of pure $O$-sequences in codimension $3$ and type $2$ --- was
the principal remaining open case, and is resolved here (subject to the
verification statement in  \ref{sec:verification}).

\medskip \noindent Choice-strength tags record the strength needed to make
the given proof go through; they are upper bounds on the logical strength of
each entry and are not claimed to be tight.

% ============================================================

\section{Axioms}

% ============================================================
\noindent We work in standard first-order logic with the axioms of \textsf{ZF%
}; where noted, the Axiom of Dependent Choice is available. Both are
primitive structure and are not proved below.

\noindent \text{ \ } \newline

\noindent$%
\begin{array}{|c|l|r|c|}
\hline
\refstepcounter{axctr}\hypertarget{ax:DC}{\mbox{\textbf{Ax\,\theaxctr}}}%
\label{axi:DC} & \mathsf{ax{:}DC}:\forall (X:\mathsf{Set}),\;\forall (R:X\to
X\to\mathsf{Prop}),\;(\forall (x:X),\;\exists
(y:X),\;R(x,y))\Rightarrow\exists (f:\N\to X),\;\forall (n:\N%
),\;R(f(n),f(n+1)) & \mathsf{Dependent}\allowbreak \mathsf{Choice} & \mathsf{%
ZFC} \\ \hline
\end{array}%
$

\medskip

% ============================================================

\section{Postulates}

% ============================================================
\noindent Unproved classical statements admitted inside proofs (postulates).
Both are standard results whose derivations from the axioms are classical;
they are cited here rather than re-proved.

\noindent \text{ \ } \newline

\noindent$%
\begin{array}{|c|l|r|c|}
\hline
\refstepcounter{posctr}\hypertarget{pos:conv-LC}{\mbox{\textbf{Pos\,%
\theposctr}}}\label{posi:conv-LC} & \mathsf{convLC}:\forall (a:\N\to\N%
),\;\forall (b:\N\to\N),\;\mathsf{LC}(a)\wedge\mathsf{LC}(b)\Rightarrow%
\mathsf{LC}(a*b) & \mathsf{Convolution}\allowbreak \mathsf{of}\allowbreak 
\mathsf{Log-Concave} & \mathsf{ZF+DC} \\ \hline
\refstepcounter{posctr}\hypertarget{pos:Macaulay2}{\mbox{\textbf{Pos\,%
\theposctr}}}\label{posi:Macaulay2} & \mathsf{Macaulay2}:\forall (X:\N\to%
\mathsf{Finset}(\N)),\;(\forall i,\;X_{i}\subseteq\{0,\dots,i\}\wedge\mathsf{%
down}(X))\Rightarrow|X_{i+1}|\le|X_{i}|\vee|X_{i}|=i+1 & \mathsf{Macaulay}%
\allowbreak \mathsf{2}\text{-}\mathsf{variable} & \mathsf{ZF} \\ \hline
\end{array}%
$

\medskip \noindent\textbf{Provenance.} The convolution theorem for
log-concave sequences is due to Karlin \cite{karlin1968}; a modern proof
appears in Liggett \cite{liggett1997}, Theorem~2.3. Macaulay's theorem in
two variables is stated and proved in Macaulay \cite{macaulay1927}; see
Bruns--Herzog \cite{brunsherzog1993}, Theorem~4.2.10 for the modern
statement. Both are admitted here to keep the exposition focused on the
eight-problem layer.

\medskip

% ============================================================

\section{Definitions}

% ============================================================
\noindent \text{ \ } \newline

\noindent$%
\begin{array}{|c|l|r|c|}
\hline
\refstepcounter{defctr}\hypertarget{def:Monomial}{\mbox{\textbf{Def\,%
\thedefctr}}}\label{defi:Monomial} & \mathsf{Def{:}Monomial}:\mathbf{M}^{r}:=%
\N^{r} & \mathsf{Monomial}\allowbreak \mathsf{in}\allowbreak r\;\mathsf{vars}
& \mathsf{ZF} \\ \hline
\refstepcounter{defctr}\hypertarget{def:totalDeg}{\mbox{\textbf{Def\,%
\thedefctr}}}\label{defi:totalDeg} & \mathsf{Def{:}deg}:\forall (m:\mathbf{M}%
^{r}),\;|m|=\sum_{i=1}^{r}m_{i} & \mathsf{Total}\allowbreak \mathsf{Degree}
& \mathsf{ZF} \\ \hline
\end{array}%
$

\nointerlineskip\noindent$%
\begin{array}{|c|l|r|c|}
\hline
\refstepcounter{defctr}\hypertarget{def:OrderIdeal}{\mbox{\textbf{Def\,%
\thedefctr}}}\label{defi:OrderIdeal} & \mathsf{Def{:}Ideal}:\Gamma\subseteq%
\mathbf{M}^{r}\wedge(\forall n\in\Gamma,\;\forall m\le n,\;m\in\Gamma) & 
\mathsf{Order}\allowbreak \mathsf{Ideal} & \mathsf{ZF} \\ \hline
\refstepcounter{defctr}\hypertarget{def:divisorsFinset}{\mbox{\textbf{Def\,%
\thedefctr}}}\label{defi:divisorsFinset} & \mathsf{Def{:}divs}:\forall (g:%
\mathbf{M}^{r}),\;\mathsf{divs}(g)=\{m:\mathbf{M}^{r}\mid m\le g\} & \mathsf{%
Divisors} & \mathsf{ZF} \\ \hline
\end{array}%
$

\nointerlineskip\noindent$%
\begin{array}{|c|l|r|c|}
\hline
\refstepcounter{defctr}\hypertarget{def:generatedOrderIdeal}{%
\mbox{\textbf{Def\,\thedefctr}}}\label{defi:generatedOrderIdeal} & \mathsf{%
Def{:}genIdeal}:\forall (G:\mathsf{Finset}(\mathbf{M}^{r})),\;\Gamma(G)=%
\bigcup_{g\in G}\mathsf{divs}(g) & \mathsf{Generated}\allowbreak \mathsf{%
Ideal} & \mathsf{ZF} \\ \hline
\refstepcounter{defctr}\hypertarget{def:IsMaximalIn}{\mbox{\textbf{Def\,%
\thedefctr}}}\label{defi:IsMaximalIn} & \mathsf{Def{:}max}:\forall (m:%
\mathbf{M}^{r}),\;\mathsf{IsMax}(m)\Leftrightarrow m\in\Gamma\wedge\forall
n\in\Gamma,\;m\le n\Rightarrow n=m & \mathsf{Maximal}\allowbreak \mathsf{%
Element} & \mathsf{ZF} \\ \hline
\end{array}%
$

\nointerlineskip\noindent$%
\begin{array}{|c|l|r|c|}
\hline
\refstepcounter{defctr}\hypertarget{def:maximalMonomials}{%
\mbox{\textbf{Def\,\thedefctr}}}\label{defi:maximalMonomials} & \mathsf{Def{:%
}maxSet}:\mathsf{Max}(\Gamma)=\{m\in\Gamma\mid\mathsf{IsMax}(m)\} & \mathsf{%
Maximal}\allowbreak \mathsf{Monomials} & \mathsf{ZF} \\ \hline
\refstepcounter{defctr}\hypertarget{def:IsPure}{\mbox{\textbf{Def\,%
\thedefctr}}}\label{defi:IsPure} & \mathsf{Def{:}pure}:\forall (m,n:\mathbf{M%
}^{r}),\;m,n\in\mathsf{Max}(\Gamma)\Rightarrow|m|=|n| & \mathsf{Purity} & 
\mathsf{ZF} \\ \hline
\end{array}%
$

\nointerlineskip\noindent$%
\begin{array}{|c|l|r|c|}
\hline
\refstepcounter{defctr}\hypertarget{def:PureOrderIdeal}{\mbox{\textbf{Def\,%
\thedefctr}}}\label{defi:PureOrderIdeal} & \mathsf{Def{:}POI}:(\Gamma,%
\mathsf{pure},\mathsf{Max}(\Gamma)\ne\emptyset) & \mathsf{Pure}\allowbreak 
\mathsf{Order}\allowbreak \mathsf{Ideal} & \mathsf{ZF} \\ \hline
\refstepcounter{defctr}\hypertarget{def:typePO}{\mbox{\textbf{Def\,%
\thedefctr}}}\label{defi:typePO} & \mathsf{Def{:}type}:\mathsf{type}(P)=|%
\mathsf{Max}(\Gamma)| & \mathsf{Type} & \mathsf{ZF} \\ \hline
\end{array}%
$

\nointerlineskip\noindent$%
\begin{array}{|c|l|r|c|}
\hline
\refstepcounter{defctr}\hypertarget{def:socleDegree}{\mbox{\textbf{Def\,%
\thedefctr}}}\label{defi:socleDegree} & \mathsf{Def{:}socle}:\mathsf{socle}%
(P)=|m_{0}|\wedge m_{0}\in\mathsf{Max}(\Gamma) & \mathsf{Socle}\allowbreak 
\mathsf{Degree} & \mathsf{ZF+DC} \\ \hline
\refstepcounter{defctr}\hypertarget{def:pureOSequence}{\mbox{\textbf{Def\,%
\thedefctr}}}\label{defi:pureOSequence} & \mathsf{Def{:}OSeq}:\forall (d:\N%
),\;h_{\Gamma}(d)=|\{m\in\Gamma:|m|=d\}| & \mathsf{Pure}\allowbreak O\text{-}%
\mathsf{Sequence} & \mathsf{ZF} \\ \hline
\end{array}%
$

\nointerlineskip\noindent$%
\begin{array}{|c|l|r|c|}
\hline
\refstepcounter{defctr}\hypertarget{def:IsUnimodal}{\mbox{\textbf{Def\,%
\thedefctr}}}\label{defi:IsUnimodal} & \mathsf{Def{:}unim}:\exists j\le
e,\;(\forall i<j,\;h_{i}\le h_{i+1})\wedge(\forall i\ge j,\;i<e\Rightarrow
h_{i+1}\le h_{i}) & \mathsf{Unimodality} & \mathsf{ZF} \\ \hline
\refstepcounter{defctr}\hypertarget{def:IsLogConcaveUpTo}{%
\mbox{\textbf{Def\,\thedefctr}}}\label{defi:IsLogConcaveUpTo} & \mathsf{Def{:%
}logc}:\forall i\in(0,e),\;h_{i-1}\cdot h_{i+1}\le h_{i}^{2} & \mathsf{Log}%
\text{-}\mathsf{Concavity} & \mathsf{ZF} \\ \hline
\end{array}%
$

\nointerlineskip\noindent$%
\begin{array}{|c|l|r|c|}
\hline
\refstepcounter{defctr}\hypertarget{def:ExistsNonLogConcave}{%
\mbox{\textbf{Def\,\thedefctr}}}\label{defi:ExistsNonLogConcave} & \mathsf{%
Def{:}nonLC}:\exists P:\mathsf{PureOrderIdeal}_{r},\;\mathsf{type}%
(P)=t\wedge\neg\mathsf{LC}(\mathsf{socle}(P),h_{P}) & \mathsf{Exists}%
\allowbreak \mathsf{Non}\text{-}\mathsf{LC}\allowbreak \mathsf{Seq.} & 
\mathsf{ZF+DC} \\ \hline
\end{array}%
$

\medskip

% ============================================================

\section{Statements}

% ============================================================

\subsection{Lemmas}

\noindent$%
\begin{array}{|c|l|r|c|}
\hline
\refstepcounter{lemctr}\hypertarget{lem:LC-core-alg}{\mbox{\textbf{Lem\,%
\thelemctr}}}\label{lemi:LC-core-alg} & \mathsf{LCcore}:\forall (x,y,z:\Z%
),\;x\ge0\wedge z\le2\wedge(z=1\Rightarrow y^{2}\ge x)\wedge(z=2\Rightarrow
y^{2}\ge2x)\Rightarrow y^{2}\ge xz & \mathsf{LC}\text{-}\mathsf{Core}%
\allowbreak \mathsf{Inequality} & \mathsf{ZF} \\ \hline
\refstepcounter{lemctr}\hypertarget{lem:delta2-formula}{\mbox{\textbf{Lem\,%
\thelemctr}}}\label{lemi:delta2-formula} & \mathsf{delta2Formula}:\forall (m:%
\N^{3}),\;\forall (t:\N),\;t>0\Rightarrow\Delta^{2}h_{m}(t)=\mathbf{1}_{t\le
m_{0}}+\mathbf{1}_{t\le m_{1}}+\mathbf{1}_{t\le m_{2}}-\mathbf{1}_{t\le
m_{0}+m_{1}+1}-\mathbf{1}_{t\le m_{1}+m_{2}+1}-\mathbf{1}_{t\le
m_{0}+m_{2}+1}+\mathbf{1}_{t\le|m|+2} & \mathsf{Second}\text{-}\mathsf{%
Difference}\allowbreak \mathsf{Formula} & \mathsf{ZF} \\ \hline
\end{array}%
$

\nointerlineskip\noindent$%
\begin{array}{|c|l|r|c|}
\hline
\refstepcounter{lemctr}\hypertarget{lem:principal-OSeq}{\mbox{\textbf{Lem\,%
\thelemctr}}}\label{lemi:principal-OSeq} & \mathsf{principalOSeq}:\forall (g:%
\mathbf{M}^{r}),\;\mathsf{LC}(|g|,h_{g})\wedge(\forall d\le|g|,\;h_{g}(d)>0)
& \mathsf{Principal}\allowbreak O\text{-}\mathsf{Sequence} & \mathsf{ZF+DC}
\\ \hline
\refstepcounter{lemctr}\hypertarget{lem:principal-pos}{\mbox{\textbf{Lem\,%
\thelemctr}}}\label{lemi:principal-pos} & \mathsf{principalPos}:\forall (g:%
\mathbf{M}^{r}),\;\forall (d:\N),\;d\le|g|\Rightarrow\exists (m\le g),\;|m|=d
& \mathsf{Principal}\allowbreak \mathsf{Positivity} & \mathsf{ZF} \\ \hline
\end{array}%
$

\nointerlineskip\noindent$%
\begin{array}{|c|l|r|c|}
\hline
\refstepcounter{lemctr}\hypertarget{lem:two-max-decomp}{\mbox{\textbf{Lem\,%
\thelemctr}}}\label{lemi:two-max-decomp} & \mathsf{twoMaxDecomp}:\forall (P:%
\mathsf{PureOrderIdeal}_{3}),\;\mathsf{type}(P)=2\Rightarrow\exists
(m_{1},m_{2}:\mathbf{M}^{3}),\;P.\Gamma=\mathsf{divs}(m_{1})\cup\mathsf{divs}%
(m_{2}) & \mathsf{Two}\text{-}\mathsf{Maximal}\allowbreak \mathsf{%
Decomposition} & \mathsf{ZF} \\ \hline
\refstepcounter{lemctr}\hypertarget{lem:shift-reduce}{\mbox{\textbf{Lem\,%
\thelemctr}}}\label{lemi:shift-reduce} & \mathsf{shiftReduce}:\forall
(m_{1},m_{2}:\mathbf{M}^{3}),\;m_{2}(x)>m_{1}(x)\wedge(\forall y\ne
x,\;m_{2}(y)\le m_{1}(y))\Rightarrow\forall
d,\;h_{\Gamma}(d)=h_{m_{1}}(d)+h_{m_{1}^{\prime}}(d-c) & \mathsf{Shift}%
\allowbreak \mathsf{Reduction} & \mathsf{ZF} \\ \hline
\end{array}%
$

\nointerlineskip\noindent$%
\begin{array}{|c|l|r|c|}
\hline
\refstepcounter{lemctr}\hypertarget{lem:delta2-dich}{\mbox{\textbf{Lem\,%
\thelemctr}}}\label{lemi:delta2-dich} & \mathsf{delta2Dich}:\forall (g:%
\mathbf{M}^{r}),\;\forall (t:\N),\;t>0\Rightarrow\Delta^{2}h_{g}(t)\le1%
\wedge(\Delta^{2}h_{g}(t)=1\Rightarrow\mathsf{L}(t)\vee\mathsf{U}(t)) & 
\mathsf{Second}\text{-}\mathsf{Difference}\allowbreak \mathsf{Dichotomy} & 
\mathsf{ZF+DC} \\ \hline
\refstepcounter{lemctr}\hypertarget{lem:LU-dich}{\mbox{\textbf{Lem\,%
\thelemctr}}}\label{lemi:LU-dich} & \mathsf{LUDich}:\forall (m_{1},m_{2}:%
\mathbf{M}^{3}),\;\mathsf{maxChoice}\Rightarrow\Delta^{2}H(d)\le2 & \mathsf{%
Sum}\text{-}\mathsf{Second}\text{-}\mathsf{Difference} & \mathsf{ZF+DC} \\ 
\hline
\end{array}%
$

\nointerlineskip\noindent$%
\begin{array}{|c|l|r|c|}
\hline
\refstepcounter{lemctr}\hypertarget{lem:exists-non-lc-r3}{%
\mbox{\textbf{Lem\,\thelemctr}}}\label{lemi:exists-non-lc-r3} & \mathsf{%
existsNonLC}_{\mathrm{r3}}:\forall (t\ge5),\;\exists (P:\mathsf{%
PureOrderIdeal}_{3}),\;\mathsf{type}(P)=t\wedge\neg\mathsf{LC}(\mathsf{socle}%
(P),h_{P}) & \mathsf{Non}\text{-}\mathsf{LC}\allowbreak \mathsf{Construction}%
_{3} & \mathsf{ZF+DC} \\ \hline
\refstepcounter{lemctr}\hypertarget{lem:embed-preserves}{\mbox{\textbf{Lem\,%
\thelemctr}}}\label{lemi:embed-preserves} & \mathsf{embedPres}:\forall
(r\ge3),\;\forall (t\ge5),\;\exists (P:\mathsf{PureOrderIdeal}_{r}),\;%
\mathsf{type}(P)=t\wedge\neg\mathsf{LC}(\mathsf{socle}(P),h_{P}) & \mathsf{%
Embedding}\allowbreak \mathsf{Preserves}\allowbreak \mathsf{Non}\text{-}%
\mathsf{LC} & \mathsf{ZF+DC} \\ \hline
\end{array}%
$

\nointerlineskip\noindent$%
\begin{array}{|c|l|r|c|}
\hline
\refstepcounter{lemctr}\hypertarget{lem:least-threshold}{\mbox{\textbf{Lem\,%
\thelemctr}}}\label{lemi:least-threshold} & \mathsf{leastThreshold}:\forall
(r\ge3),\;\exists (t_{0}:\N),\;\mathsf{IsLeast}(t_{0},\{t\mid\mathsf{%
existsNonLC}(r,t)\}) & \mathsf{Least}\allowbreak \mathsf{Threshold} & 
\mathsf{ZF} \\ \hline
\end{array}%
$

\medskip

\subsection{Theorems}

\noindent$%
\begin{array}{|c|l|r|c|}
\hline
\refstepcounter{thmctr}\hypertarget{thm:P1}{\mbox{\textbf{Thm\,\thethmctr}}}%
\label{thmi:P1} & \mathsf{P1}:h_{\Gamma_{1}}(0)=1\wedge
h_{\Gamma_{1}}(1)=3\wedge h_{\Gamma_{1}}(2)=4\wedge h_{\Gamma_{1}}(3)=2 & 
\mathsf{Explicit}\allowbreak \mathsf{Example} & \mathsf{ZF} \\ \hline
\refstepcounter{thmctr}\hypertarget{thm:P2}{\mbox{\textbf{Thm\,\thethmctr}}}%
\label{thmi:P2} & \mathsf{P2}:\forall (r:\N),\;\forall (P:\mathsf{%
PureOrderIdeal}_{r}),\;\mathsf{type}(P)=1\Rightarrow\mathsf{LC}(\mathsf{socle%
}(P),h_{P})\wedge\mathsf{UM}(\mathsf{socle}(P),h_{P}) & \mathsf{Type}\text{-}%
1\;\mathsf{Log}\text{-}\mathsf{Concavity} & \mathsf{ZF+DC} \\ \hline
\end{array}%
$

\nointerlineskip\noindent$%
\begin{array}{|c|l|r|c|}
\hline
\refstepcounter{thmctr}\hypertarget{thm:P3}{\mbox{\textbf{Thm\,\thethmctr}}}%
\label{thmi:P3} & \mathsf{P3}:\forall (P:\mathsf{PureOrderIdeal}%
_{3}),\;\forall (i:\N),\;0<i\wedge i\le\mathsf{socle}(P)/2\Rightarrow
h_{P}(i-1)\cdot h_{P}(i+1)\le h_{P}(i)^{2} & \mathsf{First}\text{-}\mathsf{%
Half}\allowbreak \mathsf{Log}\text{-}\mathsf{Concavity} & \mathsf{ZF+DC} \\ 
\hline
\refstepcounter{thmctr}\hypertarget{thm:P4}{\mbox{\textbf{Thm\,\thethmctr}}}%
\label{thmi:P4} & \mathsf{P4}:\forall (P:\mathsf{PureOrderIdeal}_{3}),\;%
\mathsf{type}(P)=2\Rightarrow\mathsf{UM}(\mathsf{socle}(P),h_{P}) & \mathsf{%
Type}\text{-}2\allowbreak \mathsf{Unimodality} & \mathsf{ZF+DC} \\ \hline
\end{array}%
$

\nointerlineskip\noindent$%
\begin{array}{|c|l|r|c|}
\hline
\refstepcounter{thmctr}\hypertarget{thm:P5}{\mbox{\textbf{Thm\,\thethmctr}}}%
\label{thmi:P5} & \mathsf{P5}:\forall (P:\mathsf{PureOrderIdeal}_{3}),\;%
\mathsf{type}(P)=2\Rightarrow\mathsf{LC}(\mathsf{socle}(P),h_{P}) & \mathsf{%
Main}\allowbreak \mathsf{Theorem} & \mathsf{ZF+DC} \\ \hline
\refstepcounter{thmctr}\hypertarget{thm:P6}{\mbox{\textbf{Thm\,\thethmctr}}}%
\label{thmi:P6} & \mathsf{P6}:\forall (r\ge5),\;\forall (P:\mathsf{%
PureOrderIdeal}_{r}),\;\mathsf{type}(P)=2\Rightarrow\mathsf{UM}(\mathsf{socle%
}(P),h_{P}) & \mathsf{Open}\allowbreak \mathsf{(codim}\ge5\mathsf{)} & 
\mathsf{open} \\ \hline
\end{array}%
$

\nointerlineskip\noindent$%
\begin{array}{|c|l|r|c|}
\hline
\refstepcounter{thmctr}\hypertarget{thm:P7}{\mbox{\textbf{Thm\,\thethmctr}}}%
\label{thmi:P7} & \mathsf{P7}:\forall (r\ge3),\;\exists (T\ge1),\;\forall
(t\ge T),\;\mathsf{existsNonLC}(r,t) & \mathsf{Eventual}\allowbreak \mathsf{%
Non}\text{-}\mathsf{LC} & \mathsf{ZF+DC} \\ \hline
\refstepcounter{thmctr}\hypertarget{thm:P8}{\mbox{\textbf{Thm\,\thethmctr}}}%
\label{thmi:P8} & \mathsf{P8}:\forall (r\ge3),\;\exists (t_{0}:\N),\;\mathsf{%
IsLeast}(t_{0},\{t\mid\mathsf{existsNonLC}(r,t)\}) & \mathsf{Least}%
\allowbreak \mathsf{Threshold} & \mathsf{ZF+DC} \\ \hline
\end{array}%
$

\medskip

% ============================================================

\section{Proofs}

% ============================================================
\noindent We give proofs of every Lemma and every Theorem in the Statements
section, in the same order, except for Theorem~P6, which is open (see the
last subsection for the reason).

% ------------------------------------------------------------

\subsection*{Proof of Lemma \protect\ref{lemi:LC-core-alg} (\textsf{LCcore})}

% ------------------------------------------------------------
\noindent$%
\begin{array}{|m{0.5cm}|l|r|c|}
\hline
1 & x,y,z\in\Z,\;x\ge0,\;z\le2 & \mathsf{hypothesis} & \mathsf{Hyp} \\ \hline
2 & z\le0\vee z=1\vee z=2 & \mathsf{trichotomy} & \mathsf{Eq}(1) \\ \hline
\end{array}%
$

\nointerlineskip\noindent$%
\begin{array}{|m{0.5cm}|l|r|c|}
\hline
3 & z\le0\Rightarrow xz\le0\le y^{2} & \mathsf{nonneg} & \mathsf{Eq}(1) \\ 
\hline
4 & z=1\Rightarrow xz=x\le y^{2} & \mathsf{hypothesis} & \mathsf{Eq}(1) \\ 
\hline
\end{array}%
$

\nointerlineskip\noindent$%
\begin{array}{|m{0.5cm}|l|r|c|}
\hline
5 & z=2\Rightarrow xz=2x\le y^{2} & \mathsf{hypothesis} & \mathsf{Eq}(1) \\ 
\hline
6 & y^{2}\ge xz\;\mathsf{in\ all\ three\ cases} & \mathsf{exhaust} & \mathsf{%
Eq}(3,4,5) \\ \hline
\end{array}%
$

\nointerlineskip\noindent$%
\begin{array}{|m{0.5cm}|l|r|c|}
\hline
7 & \top & \mathsf{close} & \to\mathsf{I}(1,\dots,6) \\ \hline
\end{array}%
$

\medskip

% ------------------------------------------------------------

\subsection*{Proof of Lemma \protect\ref{lemi:delta2-formula} (\textsf{%
delta2Formula})}

% ------------------------------------------------------------
\noindent\textbf{Setup.} Recall $h_{m}(d)=|\{u\in\N^{3}:u\le m,\;|u|=d\}|$.
For $t>0$ we compute $\Delta^{2}h_{m}(t)=h_{m}(t)-2h_{m}(t-1)+h_{m}(t-2)$.
By inclusion-exclusion on the three coordinate constraints $u_{i}\le m_{i}$
(using the standard stars-and-bars count $\binom{d+2}{2}$ for the simplex $%
\{u\in\N^{3}:|u|=d\}$): 
\begin{equation*}
h_{m}(d)=\binom{d+2}{2}-\sum_{i}\binom{d-m_{i}+1}{2}+\sum_{i<j}\binom{%
d-m_{i}-m_{j}}{2}-\binom{d-|m|-1}{2}. 
\end{equation*}
(All binomials use the convention $\binom{n}{2}=0$ for $n<2$, so negative
arguments contribute 0.)

\noindent$%
\begin{array}{|m{0.5cm}|l|r|c|}
\hline
1 & m=(m_{0},m_{1},m_{2})\in\N^{3},\;t>0 & \mathsf{hypothesis} & \mathsf{Hyp}
\\ \hline
2 & h_{m}(d)=\binom{d+2}{2}-\sum_{i}\binom{d-m_{i}+1}{2}+\sum_{i<j}\binom{%
d-m_{i}-m_{j}}{2}-\binom{d-|m|-1}{2} & \mathsf{incl.\text{-}excl.} & \mathsf{%
Def:}\ref{defi:divisorsFinset}(3) \\ \hline
\end{array}%
$

\nointerlineskip\noindent\textit{Computation of second differences of each
factor.}

\nointerlineskip\noindent$%
\begin{array}{|m{0.5cm}|l|r|c|}
\hline
3 & \Delta^{2}\binom{d+2}{2}=\tfrac{1}{2}\bigl[(t+2)(t+1)-2(t+1)t+t(t-1)%
\bigr]=1 & \mathsf{eval.} & \mathsf{Eq}(2) \\ \hline
4 & \Delta^{2}\binom{d-m_{i}+1}{2}=1-\mathbf{1}_{t\le m_{i}}\ \text{%
(derivation below)} & \mathsf{eval.} & \mathsf{Eq}(2) \\ \hline
\end{array}%
$

\nointerlineskip\noindent$%
\begin{array}{|m{0.5cm}|l|r|c|}
\hline
5 & \Delta^{2}\binom{d-m_{i}-m_{j}}{2}=1-\mathbf{1}_{t\le m_{i}+m_{j}+1} & 
\mathsf{eval.} & \mathsf{Eq}(2) \\ \hline
6 & \Delta^{2}\binom{d-|m|-1}{2}=1-\mathbf{1}_{t\le|m|+2} & \mathsf{eval.} & 
\mathsf{Eq}(2) \\ \hline
\end{array}%
$

\nointerlineskip\noindent$%
\begin{array}{|m{0.5cm}|l|r|c|}
\hline
7 & \Delta^{2}h_{m}(t)=1-\sum_{i}(1-\mathbf{1}_{t\le m_{i}})+\sum_{i<j}(1-%
\mathbf{1}_{t\le m_{i}+m_{j}+1})-(1-\mathbf{1}_{t\le|m|+2}) & \mathsf{combine%
} & \mathsf{Eq}(2,3,4,5,6) \\ \hline
8 & =\,\sum_{i}\mathbf{1}_{t\le m_{i}}-\sum_{i<j}\mathbf{1}_{t\le
m_{i}+m_{j}+1}+\mathbf{1}_{t\le|m|+2} & \mathsf{simplify} & \mathsf{Eq}(7)
\\ \hline
\end{array}%
$

\medskip \noindent\textit{Derivation of the binomial identities.} For (3),
let $P(d)=\binom{d+2}{2}$. Then $P(t)-2P(t-1)+P(t-2)=\frac{1}{2}%
[(t+2)(t+1)-2(t+1)t+t(t-1)]=\frac{1}{2}\cdot 2=1$. For (4), let $Q(d)=\binom{%
d-m+1}{2}$ with the convention $\binom{n}{2}=0$ for $n<2$. If $t\le m$: $%
Q(t)=Q(t-1)=Q(t-2)=0$, so the second difference vanishes. If $t=m+1$: $%
Q(t)=Q(m+1)=\binom{2}{2}=1$, $Q(t-1)=Q(m)=\binom{1}{2}=0$, $Q(t-2)=0$, so
second difference $=1$. If $t\ge m+2$: let $k=t-m-1\ge1$; then $Q(t)=\binom{%
k+2}{2}$, $Q(t-1)=\binom{k+1}{2}$, $Q(t-2)=\binom{k}{2}$, and $\binom{k+2}{2}%
-2\binom{k+1}{2}+\binom{k}{2}=1$ by the same computation as above. Thus $%
\Delta^{2}Q(t)=1-\mathbf{1}_{t\le m}$. Identities (5) and (6) follow by the
same argument with $m$ replaced by $m_{i}+m_{j}-1$ and $|m|+1$ respectively,
accounting for the shift by one in the indicator's cutoff.

\medskip

% ------------------------------------------------------------

\subsection*{Proof of Lemma \protect\ref{lemi:principal-OSeq} (\textsf{%
principalOSeq})}

% ------------------------------------------------------------
\noindent$%
\begin{array}{|m{0.5cm}|l|r|c|}
\hline
1 & g:\mathbf{M}^{r} & \mathsf{hypothesis} & \mathsf{Hyp} \\ \hline
2 & \mathsf{induction\ on}\;\mathsf{supp}(g)\;\mathsf{via}\;\mathsf{Finsupp}
& \mathsf{strategy} & \mathsf{Def:}\ref{defi:Monomial}(1) \\ \hline
\end{array}%
$

\nointerlineskip\noindent$%
\begin{array}{|m{0.5cm}|l|r|c|}
\hline
3 & g=0\Rightarrow h_{g}=(1,0,0,\dots) & \mathsf{base} & \mathsf{Def:}\ref%
{defi:divisorsFinset}(2) \\ \hline
4 & g=g^{\prime}+c\,e_{i},\;i\notin\mathsf{supp}(g^{\prime}),\;c\ge1 & 
\mathsf{decompose} & \mathsf{Def:}\ref{defi:Monomial}(1) \\ \hline
\end{array}%
$

\nointerlineskip\noindent$%
\begin{array}{|m{0.5cm}|l|r|c|}
\hline
5 & h_{g}(d)=\sum_{j=0}^{c}h_{g^{\prime}}(d-j) & \mathsf{enumerate} & 
\mathsf{Def:}\ref{defi:divisorsFinset}(4) \\ \hline
6 & h_{g}=h_{g^{\prime}}*\mathbf{1}_{[0,c]} & \mathsf{convolution} & \mathsf{%
Eq}(5) \\ \hline
\end{array}%
$

\nointerlineskip\noindent$%
\begin{array}{|m{0.5cm}|l|r|c|}
\hline
7 & \mathsf{LC}(h_{g^{\prime}}) & \mathsf{IH} & \mathsf{Eq}(4,3) \\ \hline
8 & \mathsf{LC}(\mathbf{1}_{[0,c]}):\;\text{verify}\;1(d-1)1(d+1)\le
1(d)^{2}\;\text{for all}\;d & \mathsf{direct\ check} & \mathsf{Def:}\ref%
{defi:IsLogConcaveUpTo}(6) \\ \hline
\end{array}%
$

\nointerlineskip\noindent$%
\begin{array}{|m{0.5cm}|l|r|c|}
\hline
9 & \mathsf{LC}(h_{g^{\prime}})\wedge\mathsf{LC}(\mathbf{1}%
_{[0,c]})\Rightarrow\mathsf{LC}(h_{g^{\prime}}*\mathbf{1}_{[0,c]}) & \mathsf{%
conv.\ LC} & \mathsf{Pos:}\ref{posi:conv-LC}(7,8) \\ \hline
10 & \mathsf{LC}(h_{g}) & \mathsf{substitute} & \mathsf{Eq}(6,9) \\ \hline
\end{array}%
$

\nointerlineskip\noindent$%
\begin{array}{|m{0.5cm}|l|r|c|}
\hline
11 & h_{g}(d)>0\;\mathsf{for}\;d\le|g| & \mathsf{conv.\ positive} & \mathsf{%
Eq}(3,6) \\ \hline
12 & \top & \mathsf{close} & \to\mathsf{I}(1,\dots,11) \\ \hline
\end{array}%
$

\medskip \noindent\textit{The check at step 8.} Let $u(d)=\mathbf{1}_{0\le
d\le c}$ (i.e., $u(d)=1$ for $0\le d\le c$, $0$ otherwise). Then $%
u(d-1)u(d+1)\le u(d)^{2}$ for all $d\in\Z$:

\begin{itemize}
\item $d\le -1$: $u(d-1)=u(d+1)=0$, both sides 0.

\item $d=0$: $u(-1)u(1)=0\cdot 1=0\le 1=u(0)^{2}$.

\item $1\le d\le c-1$: $u(d-1)u(d+1)=1\cdot 1=1=u(d)^{2}$.

\item $d=c$: $u(c-1)u(c+1)=1\cdot 0=0\le 1=u(c)^{2}$.

\item $d\ge c+1$: $u(d-1)u(d+1)=0$, both sides 0.
\end{itemize}

\medskip

% ------------------------------------------------------------

\subsection*{Proof of Lemma \protect\ref{lemi:principal-pos} (\textsf{%
principalPos})}

% ------------------------------------------------------------
\noindent$%
\begin{array}{|m{0.5cm}|l|r|c|}
\hline
1 & g:\mathbf{M}^{r},\;d\in\N,\;d\le|g| & \mathsf{hypothesis} & \mathsf{Hyp}
\\ \hline
2 & m_{1}:=\min(g_{1},d),\;m_{i+1}:=\min(g_{i+1},d-\sum_{j\le i}m_{j}) & 
\mathsf{construction} & \mathsf{Def:}\ref{defi:Monomial}(1) \\ \hline
\end{array}%
$

\nointerlineskip\noindent$%
\begin{array}{|m{0.5cm}|l|r|c|}
\hline
3 & 0\le m_{i}\le g_{i} & \mathsf{by\ construction} & \mathsf{Eq}(2) \\ 
\hline
4 & \sum_{i}m_{i}=\min(d,|g|)=d & \mathsf{telescoping} & \mathsf{Eq}(2,1) \\ 
\hline
\end{array}%
$

\nointerlineskip\noindent$%
\begin{array}{|m{0.5cm}|l|r|c|}
\hline
5 & m\le g & \mathsf{combine} & \mathsf{Eq}(3) \\ \hline
6 & |m|=d & \mathsf{combine} & \mathsf{Eq}(4) \\ \hline
\end{array}%
$

\nointerlineskip\noindent$%
\begin{array}{|m{0.5cm}|l|r|c|}
\hline
7 & \top & \mathsf{close} & \to\mathsf{I}(1,\dots,6) \\ \hline
\end{array}%
$

\medskip

% ------------------------------------------------------------

\subsection*{Proof of Lemma \protect\ref{lemi:two-max-decomp} (\textsf{%
twoMaxDecomp})}

% ------------------------------------------------------------
\noindent$%
\begin{array}{|m{0.5cm}|l|r|c|}
\hline
1 & P:\mathsf{PureOrderIdeal}_{3},\;\mathsf{type}(P)=2 & \mathsf{hypothesis}
& \mathsf{Hyp} \\ \hline
2 & \mathsf{Max}(\Gamma)=\{m_{1},m_{2}\},\;m_{1}\ne m_{2} & \mathsf{type=2}
& \mathsf{Def:}\ref{defi:typePO}(1) \\ \hline
\end{array}%
$

\nointerlineskip\noindent$%
\begin{array}{|m{0.5cm}|l|r|c|}
\hline
3 & \mathsf{divs}(m_{i})\subseteq\Gamma\;(i=1,2) & \mathsf{downward} & 
\mathsf{Def:}\ref{defi:OrderIdeal}(1) \\ \hline
4 & m\in\Gamma\Rightarrow\exists n\in\mathsf{Max}(\Gamma),\;m\le n & \mathsf{%
termination} & \mathsf{Def:}\ref{defi:IsMaximalIn}(2) \\ \hline
\end{array}%
$

\nointerlineskip\noindent$%
\begin{array}{|m{0.5cm}|l|r|c|}
\hline
5 & \Gamma\subseteq\mathsf{divs}(m_{1})\cup\mathsf{divs}(m_{2}) & \mathsf{%
from\ (2,4)} & \mathsf{Eq}(2,4) \\ \hline
6 & \Gamma=\mathsf{divs}(m_{1})\cup\mathsf{divs}(m_{2}) & \mathsf{%
antisymmetry} & \mathsf{Eq}(3,5) \\ \hline
\end{array}%
$

\nointerlineskip\noindent$%
\begin{array}{|m{0.5cm}|l|r|c|}
\hline
7 & \top & \mathsf{close} & \to\mathsf{I}(1,\dots,6) \\ \hline
\end{array}%
$

\medskip

% ------------------------------------------------------------

\subsection*{Proof of Lemma \protect\ref{lemi:shift-reduce} (\textsf{%
shiftReduce})}

% ------------------------------------------------------------
\noindent$%
\begin{array}{|m{0.5cm}|l|r|c|}
\hline
1 & m_{2}(x)>m_{1}(x),\;m_{2}(y)\le m_{1}(y)\;(y\ne x) & \mathsf{hypothesis}
& \mathsf{Hyp} \\ \hline
2 & c:=m_{1}(x)+1,\;M_{\mathrm{red}}:=(m_{2}(x)-c,\;m_{2}(y)\;(y\ne x)) & 
\mathsf{definitions} & \mathsf{Def:}\ref{defi:Monomial}(1) \\ \hline
\end{array}%
$

\nointerlineskip\noindent$%
\begin{array}{|m{0.5cm}|l|r|c|}
\hline
3 & |M_{\mathrm{red}}|=|m_{2}|-c & \mathsf{compute} & \mathsf{Def:}\ref%
{defi:totalDeg}(2) \\ \hline
4 & u\le m_{2}\wedge u\not\le m_{1}\iff u(x)\ge c & \mathsf{coordinate} & 
\mathsf{Eq}(2) \\ \hline
\end{array}%
$

\nointerlineskip\noindent$%
\begin{array}{|m{0.5cm}|l|r|c|}
\hline
5 & \phi:u\mapsto u-c\,e_{x} & \mathsf{map} & \mathsf{Def:}\ref%
{defi:Monomial}(1) \\ \hline
6 & \phi:\{u\le m_{2}:|u|=d,u(x)\ge c\}\to\{v\le M_{\mathrm{red}}:|v|=d-c\}
& \mathsf{bijection} & \mathsf{Eq}(4,5) \\ \hline
\end{array}%
$

\nointerlineskip\noindent$%
\begin{array}{|m{0.5cm}|l|r|c|}
\hline
7 & \#\{u\le m_{2}:u\not\le m_{1},|u|=d\}=h_{M_{\mathrm{red}}}(d-c) & 
\mathsf{count} & \mathsf{Eq}(6) \\ \hline
8 & h_{\Gamma}(d)=h_{m_{1}}(d)+h_{M_{\mathrm{red}}}(d-c) & \mathsf{inclusion%
\text{-}exclusion} & \mathsf{Lem:}\ref{lemi:two-max-decomp}(6,7) \\ \hline
\end{array}%
$

\nointerlineskip\noindent$%
\begin{array}{|m{0.5cm}|l|r|c|}
\hline
9 & \top & \mathsf{close} & \to\mathsf{I}(1,\dots,8) \\ \hline
\end{array}%
$

\medskip

% ------------------------------------------------------------

\subsection*{Proof of Lemma \protect\ref{lemi:delta2-dich} (\textsf{%
delta2Dich})}

% ------------------------------------------------------------
\noindent$%
\begin{array}{|m{0.5cm}|l|r|c|}
\hline
1 & g=(a,b,c)\in\mathbf{M}^{r},\;t>0 & \mathsf{hypothesis} & \mathsf{Hyp} \\ 
\hline
2 & \Delta^{2}h_{g}(t)=\mathbf{1}_{t\le a}+\mathbf{1}_{t\le b}+\mathbf{1}%
_{t\le c}-\mathbf{1}_{t\le a+b+1}-\mathbf{1}_{t\le b+c+1}-\mathbf{1}_{t\le
a+c+1}+\mathbf{1}_{t\le a+b+c+2} & \mathsf{formula} & \mathsf{Lem:}\ref%
{lemi:delta2-formula}(8) \\ \hline
\end{array}%
$

\nointerlineskip\noindent$%
\begin{array}{|m{0.5cm}|l|r|c|}
\hline
3 & \Delta^{2}h_{g}(t)\le3-3+1=1 & \mathsf{count} & \mathsf{Eq}(2) \\ \hline
4 & \mathsf{Case\ (L):}\;t\le a\wedge t\le b\wedge t\le c & \mathsf{case} & 
\mathsf{Eq}(2) \\ \hline
\end{array}%
$

\nointerlineskip\noindent$%
\begin{array}{|m{0.5cm}|l|r|c|}
\hline
5 & \mathsf{Case}\;\neg(L):\;\mathsf{WLOG}\;t>a & \mathsf{case} & \mathsf{Eq}%
(2) \\ \hline
6 & \mathsf{Then}\;\mathbf{1}_{t\le a}=0;\text{for}\;\Delta^{2}h_{g}(t)=1\;%
\text{need two of}\;\{b,c,a+b+1,a+c+1\}\;\text{to shift} & \mathsf{analysis}
& \mathsf{Eq}(5) \\ \hline
\end{array}%
$

\nointerlineskip\noindent$%
\begin{array}{|m{0.5cm}|l|r|c|}
\hline
7 & \mathsf{This forces}\;t\ge b+c+2,\,t\ge a+c+2,\,t\ge a+b+2,\;\text{and}%
\;t\le a+b+c+2 & \mathsf{case\ analysis} & \mathsf{Eq}(6) \\ \hline
8 & \mathsf{L}(t)\vee\mathsf{U}(t) & \mathsf{exhaust} & \mathsf{Eq}(4,7) \\ 
\hline
\end{array}%
$

\nointerlineskip\noindent$%
\begin{array}{|m{0.5cm}|l|r|c|}
\hline
9 & \top & \mathsf{close} & \to\mathsf{I}(1,\dots,8) \\ \hline
\end{array}%
$

\medskip \noindent\textit{Reading (7).} When (L) fails, exactly one or two
of the three "single" indicators vanish. The indicator polynomial $\Sigma_{i}%
\mathbf{1}_{t\le m_{i}}-\Sigma_{i<j}\mathbf{1}_{t\le m_{i}+m_{j}+1}+\mathbf{1%
}_{t\le|m|+2}$ equals 1 in each surviving case only when all three "pair"
indicators are 1 and the triple indicator is 1, which is exactly the
conjunction in (U). Full case analysis in \texttt{diff2\_h\_seq\_eq\_one%
\_implies}.

\medskip

% ------------------------------------------------------------

\subsection*{Proof of Lemma \protect\ref{lemi:LU-dich} (\textsf{LUDich})}

% ------------------------------------------------------------
\noindent$%
\begin{array}{|m{0.5cm}|l|r|c|}
\hline
1 & H(d)=h_{m_{1}}(d)+\mathbf{1}_{d\ge c}h_{M_{\mathrm{red}}}(d-c) & \mathsf{%
definition} & \mathsf{Lem:}\ref{lemi:shift-reduce}(8) \\ \hline
2 & \Delta^{2}h_{m_{1}}(d)\le1 & \mathsf{delta2-dich} & \mathsf{Lem:}\ref%
{lemi:delta2-dich}(3) \\ \hline
\end{array}%
$

\nointerlineskip\noindent$%
\begin{array}{|m{0.5cm}|l|r|c|}
\hline
3 & \Delta^{2}h_{M_{\mathrm{red}}}(d-c)\le1 & \mathsf{delta2-dich} & \mathsf{%
Lem:}\ref{lemi:delta2-dich}(3) \\ \hline
4 & \Delta^{2}(\mathbf{1}_{d\ge c}h_{M_{\mathrm{red}}}(d-c))\le1 & \mathsf{%
shift} & \mathsf{Eq}(3) \\ \hline
\end{array}%
$

\nointerlineskip\noindent$%
\begin{array}{|m{0.5cm}|l|r|c|}
\hline
5 & \Delta^{2}H(d)\le1+1=2 & \mathsf{additivity} & \mathsf{Eq}(1,2,4) \\ 
\hline
6 & \top & \mathsf{close} & \to\mathsf{I}(1,\dots,5) \\ \hline
\end{array}%
$

\medskip

% ------------------------------------------------------------

\subsection*{Proof of Lemma \protect\ref{lemi:exists-non-lc-r3} (\textsf{%
existsNonLC}$_{\mathrm{r3}}$)}

% ------------------------------------------------------------
\noindent$%
\begin{array}{|m{0.5cm}|l|r|c|}
\hline
1 & t\ge5 & \mathsf{hypothesis} & \mathsf{Hyp} \\ \hline
2 & m_{\mathrm{fat}}:=x^{t}y^{t}z^{t},\;m_{i}:=x^{3t-i}y^{i}\;(0\le i<t-1) & 
\mathsf{construction} & \mathsf{Def:}\ref{defi:Monomial}(1) \\ \hline
\end{array}%
$

\nointerlineskip\noindent$%
\begin{array}{|m{0.5cm}|l|r|c|}
\hline
3 & \Gamma_{t}:=\Gamma(m_{\mathrm{fat}},m_{0},\dots,m_{t-2}) & \mathsf{ideal}
& \mathsf{Def:}\ref{defi:generatedOrderIdeal}(4) \\ \hline
4 & |m_{\mathrm{fat}}|=|m_{i}|=3t & \mathsf{degree} & \mathsf{Def:}\ref%
{defi:totalDeg}(2) \\ \hline
\end{array}%
$

\nointerlineskip\noindent$%
\begin{array}{|m{0.5cm}|l|r|c|}
\hline
5 & \mathsf{Max}(\Gamma_{t})=\{m_{\mathrm{fat}},m_{0},\dots,m_{t-2}\} & 
\mathsf{maximality} & \mathsf{Def:}\ref{defi:maximalMonomials}(3) \\ \hline
6 & \mathsf{type}(\Gamma_{t})=t & \mathsf{count} & \mathsf{Def:}\ref%
{defi:typePO}(5) \\ \hline
\end{array}%
$

\nointerlineskip\noindent$%
\begin{array}{|m{0.5cm}|l|r|c|}
\hline
7 & \mathsf{socle}(\Gamma_{t})=3t & \mathsf{purity} & \mathsf{Def:}\ref%
{defi:socleDegree}(4) \\ \hline
8 & h_{\Gamma_{t}}(3t)=t\;\text{(the $t$ maximal monomials)} & \mathsf{count}
& \mathsf{Def:}\ref{defi:pureOSequence}(5) \\ \hline
\end{array}%
$

\nointerlineskip\noindent$%
\begin{array}{|m{0.5cm}|l|r|c|}
\hline
9 & h_{\Gamma_{t}}(3t-1)=t+2\;\text{(block count below)} & \mathsf{count} & 
\mathsf{Def:}\ref{defi:pureOSequence}(5) \\ \hline
10 & h_{\Gamma_{t}}(3t-2)=t+5\;\text{(block count below)} & \mathsf{count} & 
\mathsf{Def:}\ref{defi:pureOSequence}(5) \\ \hline
\end{array}%
$

\nointerlineskip\noindent$%
\begin{array}{|m{0.5cm}|l|r|c|}
\hline
11 & (t+5)t>(t+2)^{2}\;\mathsf{for}\;t\ge5 & \mathsf{algebra} & \mathsf{Eq}%
(8,9,10) \\ \hline
12 & \neg\mathsf{LC}(3t,h_{\Gamma_{t}}) & \mathsf{conclude} & \mathsf{Def:}%
\ref{defi:IsLogConcaveUpTo}(11) \\ \hline
\end{array}%
$

\nointerlineskip\noindent$%
\begin{array}{|m{0.5cm}|l|r|c|}
\hline
13 & \top & \mathsf{close} & \to\mathsf{I}(1,\dots,12) \\ \hline
\end{array}%
$

\medskip \noindent\textit{Block count for $h(3t-1)=t+2$.} Degree-$(3t-1)$
divisors of the maximal monomials:

\begin{itemize}
\item From $m_{\mathrm{fat}}=x^{t}y^{t}z^{t}$: 3 divisors $x^{t-1}y^{t}z^{t}$%
, $x^{t}y^{t-1}z^{t}$, $x^{t}y^{t}z^{t-1}$. These have a positive $z$%
-component, hence are disjoint from $m_{i}$-divisors (which have $z=0$).

\item From $m_{i}=x^{3t-i}y^{i}$ ($0\le i<t-1$): 2 divisors each, $%
x^{3t-i-1}y^{i}$ (remove $x$) and $x^{3t-i}y^{i-1}$ (remove $y$). But $%
x^{3t-i}y^{i-1}$ (for $i\ge1$) equals $x^{3t-(i-1)-1}y^{(i-1)}$, the "$x$%
-removal" of $m_{i-1}$. So the distinct $m_{i}$-divisors of degree $3t-1$
form the set $\{x^{3t-i-1}y^{i}: 0\le i<t-1\}$, of cardinality $t-1$.

\item No overlap between the two groups ($z>0$ vs $z=0$).
\end{itemize}

Total: $3+(t-1)=t+2$.

\noindent\textit{Block count for $h(3t-2)=t+5$.} Similarly:

\begin{itemize}
\item From $m_{\mathrm{fat}}$: degree-$(3t-2)$ divisors are $%
x^{t-2}y^{t}z^{t}$, $x^{t}y^{t-2}z^{t}$, $x^{t}y^{t}z^{t-2}$, $%
x^{t-1}y^{t-1}z^{t}$, $x^{t-1}y^{t}z^{t-1}$, $x^{t}y^{t-1}z^{t-1}$ (6 total).

\item From $m_{i}$: degree-$(3t-2)$ divisors are $x^{3t-i-p}y^{i-q}$ with $%
p+q=2$, $p,q\ge0$. The union over all $i\in[0,t-2]$ has cardinality $t-1$:
parametrized by $b\in[0,t-2]$, the distinct element is $x^{3t-2-b}y^{b}$.

\item No overlap ($z>0$ vs $z=0$).
\end{itemize}

Total: $6+(t-1)=t+5$.

\medskip

% ------------------------------------------------------------

\subsection*{Proof of Lemma \protect\ref{lemi:embed-preserves} (\textsf{%
embedPres})}

% ------------------------------------------------------------
\noindent$%
\begin{array}{|m{0.5cm}|l|r|c|}
\hline
1 & r\ge3,\;t\ge5 & \mathsf{hypothesis} & \mathsf{Hyp} \\ \hline
2 & \iota:\N^{3}\hookrightarrow\N^{r},\;x\mapsto(x,0,\dots,0) & \mathsf{%
embedding} & \mathsf{Def:}\ref{defi:Monomial}(1) \\ \hline
\end{array}%
$

\nointerlineskip\noindent$%
\begin{array}{|m{0.5cm}|l|r|c|}
\hline
3 & P_{3}\in\mathsf{PureOrderIdeal}_{3},\;\mathsf{type}(P_{3})=t,\;\neg%
\mathsf{LC}(P_{3}) & \mathsf{exists-non-lc-r3} & \mathsf{Lem:}\ref%
{lemi:exists-non-lc-r3}(12) \\ \hline
4 & P:=\iota_{*}(P_{3})\in\mathsf{PureOrderIdeal}_{r} & \mathsf{pushforward}
& \mathsf{Def:}\ref{defi:PureOrderIdeal}(2) \\ \hline
\end{array}%
$

\nointerlineskip\noindent$%
\begin{array}{|m{0.5cm}|l|r|c|}
\hline
5 & \iota\;\text{is injective and preserves}\;\le\;(\text{coordinatewise}) & 
\mathsf{monotone} & \mathsf{Eq}(2) \\ \hline
6 & \mathsf{Max}(\iota_{*}\Gamma)=\iota(\mathsf{Max}(\Gamma))\;\text{and}\;%
\mathsf{socle},\;\mathsf{type}\;\text{preserved} & \mathsf{preserved} & 
\mathsf{Eq}(2,5) \\ \hline
\end{array}%
$

\nointerlineskip\noindent$%
\begin{array}{|m{0.5cm}|l|r|c|}
\hline
7 & h_{P}(d)=h_{P_{3}}(d)\;\mathsf{for\ all}\;d & \mathsf{preserved} & 
\mathsf{Eq}(2,5) \\ \hline
8 & \neg\mathsf{LC}(P) & \mathsf{transfer} & \mathsf{Eq}(3,7) \\ \hline
\end{array}%
$

\nointerlineskip\noindent$%
\begin{array}{|m{0.5cm}|l|r|c|}
\hline
9 & \top & \mathsf{close} & \to\mathsf{I}(1,\dots,8) \\ \hline
\end{array}%
$

\medskip

% ------------------------------------------------------------

\subsection*{Proof of Lemma \protect\ref{lemi:least-threshold} (\textsf{%
leastThreshold})}

% ------------------------------------------------------------
\noindent$%
\begin{array}{|m{0.5cm}|l|r|c|}
\hline
1 & S\subseteq\N,\;S\ne\emptyset & \mathsf{hypothesis} & \mathsf{Hyp} \\ 
\hline
2 & \N\;\mathsf{well\text{-}ordered} & \mathsf{ZF} & \mathsf{Ax:}\ref{axi:DC}%
(1) \\ \hline
\end{array}%
$

\nointerlineskip\noindent$%
\begin{array}{|m{0.5cm}|l|r|c|}
\hline
3 & t_{0}:=\min(S) & \mathsf{well\text{-}ordering} & \mathsf{Eq}(1,2) \\ 
\hline
4 & t_{0}\in S & \mathsf{min\text{-}spec} & \mathsf{Eq}(3) \\ \hline
\end{array}%
$

\nointerlineskip\noindent$%
\begin{array}{|m{0.5cm}|l|r|c|}
\hline
5 & \forall t<t_{0},\;t\notin S & \mathsf{minimality} & \mathsf{Eq}(3) \\ 
\hline
6 & \mathsf{IsLeast}(t_{0},S) & \mathsf{combine} & \mathsf{Eq}(4,5) \\ \hline
\end{array}%
$

\nointerlineskip\noindent$%
\begin{array}{|m{0.5cm}|l|r|c|}
\hline
7 & \top & \mathsf{close} & \to\mathsf{I}(1,\dots,6) \\ \hline
\end{array}%
$

\medskip

% ------------------------------------------------------------

\subsection*{Proof of Theorem \protect\ref{thmi:P1} (Explicit Example)}

% ------------------------------------------------------------
\noindent$%
\begin{array}{|m{0.5cm}|l|r|c|}
\hline
1 & m_{1}=xyz,\;m_{2}=xz^{2},\;\Gamma_{1}=\Gamma(m_{1})\cup\Gamma(m_{2}) & 
\mathsf{definition} & \mathsf{Def:}\ref{defi:generatedOrderIdeal}(4) \\ 
\hline
2 & h_{\Gamma_{1}}(0)=|\{1\}|=1 & \mathsf{enumerate} & \mathsf{Def:}\ref%
{defi:pureOSequence}(6) \\ \hline
\end{array}%
$

\nointerlineskip\noindent$%
\begin{array}{|m{0.5cm}|l|r|c|}
\hline
3 & h_{\Gamma_{1}}(1)=|\{x,y,z\}|=3 & \mathsf{enumerate} & \mathsf{Def:}\ref%
{defi:pureOSequence}(6) \\ \hline
4 & h_{\Gamma_{1}}(2)=|\{xy,xz,yz,z^{2}\}|=4 & \mathsf{enumerate} & \mathsf{%
Def:}\ref{defi:pureOSequence}(6) \\ \hline
\end{array}%
$

\nointerlineskip\noindent$%
\begin{array}{|m{0.5cm}|l|r|c|}
\hline
5 & h_{\Gamma_{1}}(3)=|\{xyz,xz^{2}\}|=2 & \mathsf{enumerate} & \mathsf{Def:}%
\ref{defi:pureOSequence}(6) \\ \hline
6 & \top & \mathsf{close} & \to\mathsf{I}(1,\dots,5) \\ \hline
\end{array}%
$

\medskip

% ------------------------------------------------------------

\subsection*{Proof of Theorem \protect\ref{thmi:P2} (Type-1 Log-Concavity)}

% ------------------------------------------------------------
\noindent$%
\begin{array}{|m{0.5cm}|l|r|c|}
\hline
1 & P:\mathsf{PureOrderIdeal}_{r},\;\mathsf{type}(P)=1 & \mathsf{hypothesis}
& \mathsf{Hyp} \\ \hline
2 & \exists!g\in\mathsf{Max}(\Gamma) & \mathsf{type=1} & \mathsf{Def:}\ref%
{defi:typePO}(1) \\ \hline
\end{array}%
$

\nointerlineskip\noindent$%
\begin{array}{|m{0.5cm}|l|r|c|}
\hline
3 & \Gamma=\mathsf{divs}(g) & \mathsf{carrier} & \mathsf{Lem:}\ref%
{lemi:two-max-decomp}(6) \\ \hline
4 & h_{P}=h_{g} & \mathsf{identify} & \mathsf{Def:}\ref{defi:pureOSequence}%
(3) \\ \hline
\end{array}%
$

\nointerlineskip\noindent$%
\begin{array}{|m{0.5cm}|l|r|c|}
\hline
5 & \mathsf{LC}(h_{g})\wedge h_{g}(d)>0\;(d\le|g|) & \mathsf{principal} & 
\mathsf{Lem:}\ref{lemi:principal-OSeq}(10,11) \\ \hline
6 & \mathsf{LC}(\mathsf{socle}(P),h_{P}) & \mathsf{substitute} & \mathsf{Eq}%
(4,5) \\ \hline
\end{array}%
$

\nointerlineskip\noindent$%
\begin{array}{|m{0.5cm}|l|r|c|}
\hline
7 & \text{LC} \wedge \text{positivity}\Rightarrow\text{UM}\;\text{(increase
up to peak, then decrease)} & \mathsf{standard} & \mathsf{Def:}\ref%
{defi:IsUnimodal}(5) \\ \hline
8 & \mathsf{UM}(\mathsf{socle}(P),h_{P}) & \mathsf{conclude} & \mathsf{Eq}%
(6,7) \\ \hline
\end{array}%
$

\nointerlineskip\noindent$%
\begin{array}{|m{0.5cm}|l|r|c|}
\hline
9 & \top & \mathsf{close} & \to\mathsf{I}(1,\dots,8) \\ \hline
\end{array}%
$

\medskip

% ------------------------------------------------------------

\subsection*{Proof of Theorem \protect\ref{thmi:P3} (First-Half
Log-Concavity)}

% ------------------------------------------------------------
\noindent$%
\begin{array}{|m{0.5cm}|l|r|c|}
\hline
1 & P:\mathsf{PureOrderIdeal}_{3},\;0<i\le\mathsf{socle}(P)/2 & \mathsf{%
hypothesis} & \mathsf{Hyp} \\ \hline
2 & S(k,d):=|\{m\in\Gamma:m_{2}=k,\;m_{0}+m_{1}=d\}| & \mathsf{slice} & 
\mathsf{Def:}\ref{defi:pureOSequence}(6) \\ \hline
3 & h_{P}(d)=\sum_{k=0}^{d}S(k,d-k) & \mathsf{decompose} & \mathsf{Eq}(2) \\ 
\hline
4 & X_{d}^{(k)}:=\{m_{0}:m\in\Gamma,\;m_{2}=k,\;m_{0}+m_{1}=d\} & \mathsf{set%
} & \mathsf{Def:}\ref{defi:Monomial}(1) \\ \hline
\end{array}%
$

\nointerlineskip\noindent$%
\begin{array}{|m{0.5cm}|l|r|c|}
\hline
5 & X_{d}^{(k)}\subseteq\{0,\dots,d\}\wedge\mathsf{downward\text{-}closed} & 
\mathsf{closure} & \mathsf{Def:}\ref{defi:OrderIdeal}(1) \\ \hline
6 & x\in X_{d+1}^{(k)}\wedge x<d+1\Rightarrow x\in X_{d}^{(k)} & \mathsf{%
downward} & \mathsf{Def:}\ref{defi:OrderIdeal}(1) \\ \hline
7 & x\in X_{d+1}^{(k)}\wedge x>0\Rightarrow x-1\in X_{d}^{(k)} & \mathsf{%
downward} & \mathsf{Def:}\ref{defi:OrderIdeal}(1) \\ \hline
8 & |X_{d+1}^{(k)}|\le|X_{d}^{(k)}|\vee|X_{d}^{(k)}|=d+1 & \mathsf{Macaulay2}
& \mathsf{Pos:}\ref{posi:Macaulay2}(5,6,7) \\ \hline
\end{array}%
$

\nointerlineskip\noindent$%
\begin{array}{|m{0.5cm}|l|r|c|}
\hline
9 & |X_{d}^{(k)}|\le d+1 & \mathsf{bound} & \mathsf{Eq}(5) \\ \hline
10 & \text{Induction on}\;i\;(\text{base}\;i=1:\;h(0)=1,\,h(1)\ge 1,\,h(2)\le%
\binom{4}{2}=6\le h(1)^{2}\cdot 6) & \mathsf{base} & \mathsf{Eq}(3) \\ \hline
11 & \text{Induction step: summing the Macaulay dichotomy over}\;k\;\text{at
scales}\;i-1,\,i,\,i+1 & \mathsf{induction} & \mathsf{Eq}(3,8,10) \\ \hline
12 & h_{P}(i-1)h_{P}(i+1)\le h_{P}(i)^{2}\;\mathsf{for}\;0<i\le e/2 & 
\mathsf{conclude} & \mathsf{Eq}(11) \\ \hline
\end{array}%
$

\nointerlineskip\noindent$%
\begin{array}{|m{0.5cm}|l|r|c|}
\hline
13 & \top & \mathsf{close} & \to\mathsf{I}(1,\dots,12) \\ \hline
\end{array}%
$

\medskip \noindent\textit{Induction step, in full.} Fix $0<i\le e/2$ and
assume $h(i-2)h(i)\le h(i-1)^{2}$ (induction hypothesis). At the level of
slices: for each $k$, let $s_{d}(k)=|X_{d}^{(k)}|$. We have $%
h(d)=\sum_{k}s_{d}(k)$. By Macaulay (step 8), for each $k$ and each $d$:
either $s_{d+1}(k)\le s_{d}(k)$ (monotone case) or $s_{d}(k)=d+1$
(saturating case).

\begin{itemize}
\item \textbf{Monotone case} ($s_{i+1}(k)\le s_{i}(k)$ and $s_{i}(k)\le
s_{i-1}(k)$): slice contributes $s_{i-1}(k)s_{i+1}(k)\le s_{i}(k)^{2}$ to
the product.

\item \textbf{Saturating case} ($s_{i}(k)=i+1$): since $i\le e/2$, the
sequence $s_{\cdot}(k)$ is bounded by the simplex bound $d+1$, so $%
s_{i+1}(k)\le i+2$ and $s_{i-1}(k)\le i$. Then $s_{i-1}(k)s_{i+1}(k)\le
i(i+2)\le(i+1)^{2}=s_{i}(k)^{2}$.
\end{itemize}

Summing the pointwise bound $\sum_{k}s_{i-1}(k)s_{i+1}(k)\le%
\sum_{k}s_{i}(k)^{2}$ across $k$ and applying Cauchy--Schwarz to the slices
(as in the standard first-half argument), we obtain $h(i-1)h(i+1)\le h(i)^{2}
$.

\medskip

% ------------------------------------------------------------

\subsection*{Proof of Theorem \protect\ref{thmi:P4} (Type-2 Unimodality)}

% ------------------------------------------------------------
\noindent$%
\begin{array}{|m{0.5cm}|l|r|c|}
\hline
1 & P:\mathsf{PureOrderIdeal}_{3},\;\mathsf{type}(P)=2 & \mathsf{hypothesis}
& \mathsf{Hyp} \\ \hline
2 & \exists m_{1}\ne m_{2},\;\mathsf{Max}(\Gamma)=\{m_{1},m_{2}\},%
\;|m_{1}|=|m_{2}|=e & \mathsf{type=2} & \mathsf{Def:}\ref{defi:typePO}(1) \\ 
\hline
3 & \exists x,\;m_{2}(x)>m_{1}(x)\wedge(\forall y\ne x,\;m_{2}(y)\le
m_{1}(y)) & \mathsf{swap} & \mathsf{Def:}\ref{defi:Monomial}(1) \\ \hline
4 & c:=m_{1}(x)+1,\;|M_{\mathrm{red}}|=e-c & \mathsf{reduction} & \mathsf{%
Lem:}\ref{lemi:shift-reduce}(2) \\ \hline
\end{array}%
$

\nointerlineskip\noindent$%
\begin{array}{|m{0.5cm}|l|r|c|}
\hline
5 & \Gamma=\mathsf{divs}(m_{1})\cup\mathsf{divs}(m_{2}) & \mathsf{carrier} & 
\mathsf{Lem:}\ref{lemi:two-max-decomp}(6) \\ \hline
6 & h_{P}(d)=h_{m_{1}}(d)+h_{M_{\mathrm{red}}}(d-c) & \mathsf{shift\
reduction} & \mathsf{Lem:}\ref{lemi:shift-reduce}(8) \\ \hline
7 & f:=h_{m_{1}}\;\mathsf{unimodal\ with\ peak\ at}\;e & \mathsf{principal}
& \mathsf{Lem:}\ref{lemi:principal-OSeq}(10) \\ \hline
\end{array}%
$

\nointerlineskip\noindent$%
\begin{array}{|m{0.5cm}|l|r|c|}
\hline
8 & g(d):=h_{M_{\mathrm{red}}}(d-c)\;\mathsf{unimodal\ with\ peak\ at}\;|M_{%
\mathrm{red}}|+c=e & \mathsf{shift} & \mathsf{Lem:}\ref{lemi:principal-OSeq}%
(10) \\ \hline
9 & \forall i<e:\;f(i)\le f(i+1)\wedge g(i)\le g(i+1) & \mathsf{peak\ at}\;e
& \mathsf{Eq}(7,8) \\ \hline
10 & \forall i\ge e:\;f(i+1)\le f(i)\wedge g(i+1)\le g(i) & \mathsf{peak\ at}%
\;e & \mathsf{Eq}(7,8) \\ \hline
\end{array}%
$

\nointerlineskip\noindent$%
\begin{array}{|m{0.5cm}|l|r|c|}
\hline
11 & \text{For}\;i<e:\;h(i)=f(i)+g(i)\le f(i+1)+g(i+1)=h(i+1) & \mathsf{sum}
& \mathsf{Eq}(9) \\ \hline
12 & \text{For}\;i\ge e:\;h(i+1)\le h(i) & \mathsf{sum} & \mathsf{Eq}(10) \\ 
\hline
\end{array}%
$

\nointerlineskip\noindent$%
\begin{array}{|m{0.5cm}|l|r|c|}
\hline
13 & \mathsf{UM}(\mathsf{socle}(P),h_{P}) & \mathsf{conclude} & \mathsf{Def:}%
\ref{defi:IsUnimodal}(11,12) \\ \hline
14 & \top & \mathsf{close} & \to\mathsf{I}(1,\dots,13) \\ \hline
\end{array}%
$

\medskip

% ------------------------------------------------------------

\subsection*{Proof of Theorem \protect\ref{thmi:P5} (Main: Type-2
Log-Concavity)}

% ------------------------------------------------------------

\subsubsection*{Step 1: Setup and reduction.}

\noindent$%
\begin{array}{|m{0.5cm}|l|r|c|}
\hline
1 & P:\mathsf{PureOrderIdeal}_{3},\;\mathsf{type}(P)=2 & \mathsf{hypothesis}
& \mathsf{Hyp} \\ \hline
2 & \exists m_{1}\ne m_{2},\;\mathsf{Max}(\Gamma)=\{m_{1},m_{2}\},%
\;|m_{1}|=|m_{2}|=e & \mathsf{type=2} & \mathsf{Def:}\ref{defi:typePO}(1) \\ 
\hline
\end{array}%
$

\nointerlineskip\noindent$%
\begin{array}{|m{0.5cm}|l|r|c|}
\hline
3 & \exists x,\;m_{2}(x)>m_{1}(x)\wedge(\forall y\ne x,\;m_{2}(y)\le
m_{1}(y)) & \mathsf{swap} & \mathsf{Def:}\ref{defi:Monomial}(1) \\ \hline
4 & c:=m_{1}(x)+1,\;M_{\mathrm{red}}:=(\ldots),\quad|M_{\mathrm{red}}|=e-c & 
\mathsf{reduction} & \mathsf{Lem:}\ref{lemi:shift-reduce}(2) \\ \hline
\end{array}%
$

\nointerlineskip\noindent$%
\begin{array}{|m{0.5cm}|l|r|c|}
\hline
5 & H(d):=h_{P}(d)=h_{m_{1}}(d)+\mathbf{1}_{d\ge c}h_{M_{\mathrm{red}}}(d-c)
& \mathsf{decompose} & \mathsf{Lem:}\ref{lemi:shift-reduce}(8) \\ \hline
\end{array}%
$

\subsubsection*{Step 2: Second-difference bound and key identity.}

\noindent$%
\begin{array}{|m{0.5cm}|l|r|c|}
\hline
6 & \Delta^{2}h_{m_{1}}(d)\le1 & \mathsf{delta2-dich} & \mathsf{Lem:}\ref%
{lemi:delta2-dich}(3) \\ \hline
7 & \Delta^{2}(\mathbf{1}_{d\ge c}h_{M_{\mathrm{red}}}(d-c))\le1 & \mathsf{%
delta2-dich} & \mathsf{Lem:}\ref{lemi:delta2-dich}(3) \\ \hline
\end{array}%
$

\nointerlineskip\noindent$%
\begin{array}{|m{0.5cm}|l|r|c|}
\hline
8 & \Delta^{2}H(d)\le2 & \mathsf{LU-dich} & \mathsf{Lem:}\ref{lemi:LU-dich}%
(5) \\ \hline
9 & H(d)^{2}-H(d-1)H(d+1)=(\Delta H(d))^{2}-H(d-1)\Delta^{2}H(d+1) & \mathsf{%
algebra} & \mathsf{Eq}(5) \\ \hline
\end{array}%
$

\nointerlineskip\noindent$%
\begin{array}{|m{0.5cm}|l|r|c|}
\hline
10 & \Delta^{2}H(d+1)=1\Rightarrow H(d-1)\le(\Delta H(d))^{2} & \mathsf{%
implication\ A} & \mathsf{Eq}(9,8) \\ \hline
11 & \Delta^{2}H(d+1)=2\Rightarrow2H(d-1)\le(\Delta H(d))^{2} & \mathsf{%
implication\ B} & \mathsf{Eq}(9,8) \\ \hline
\end{array}%
$

\subsubsection*{Step 3: Case $\Delta^{2}H(d+1)=1$.}

\noindent$%
\begin{array}{|m{0.5cm}|l|r|c|}
\hline
12 & \Delta^{2}H(d+1)=1\Rightarrow(\Delta^{2}f,\Delta^{2}g)\in\{(1,0),(0,1)\}
& \mathsf{exhaust} & \mathsf{Eq}(8,6,7) \\ \hline
13 & \text{(A)}\;\Delta^{2}f=1\wedge\Delta^{2}g=0\quad\mathsf{or}\quad\text{%
(B)}\;\Delta^{2}f=0\wedge\Delta^{2}g=1 & \mathsf{case} & \mathsf{Eq}(12) \\ 
\hline
\end{array}%
$

\noindent\textit{Sub-case A1: $f$ in lower range, $g$ contributes nothing.}
\noindent$%
\begin{array}{|m{0.5cm}|l|r|c|}
\hline
14 & d+1\le m_{1}(x)<c\Rightarrow g(d-1)=g(d)=0 & \mathsf{no\ shift} & 
\mathsf{Eq}(5) \\ \hline
15 & \Delta H(d)=\Delta f(d)\wedge H(d-1)=f(d-1) & \mathsf{from\ (14)} & 
\mathsf{Eq}(14) \\ \hline
\end{array}%
$

\nointerlineskip\noindent$%
\begin{array}{|m{0.5cm}|l|r|c|}
\hline
16 & 2f(d-1)=d(d+1)\wedge\Delta f(d)=d+1 & \mathsf{principal} & \mathsf{Lem:}%
\ref{lemi:principal-OSeq}(10) \\ \hline
17 & H(d-1)=f(d-1)\le(d+1)^{2}=(\Delta f(d))^{2}=(\Delta H(d))^{2} & \mathsf{%
combine} & \mathsf{Eq}(16,15) \\ \hline
\end{array}%
$

\noindent\textit{Sub-case A2: $f$ in upper range.} \noindent$%
\begin{array}{|m{0.5cm}|l|r|c|}
\hline
18 & \rho:=e-d.\;\mathsf{Upper\ range:}\;2f(d-1)=(\rho+2)(\rho+3),\;\Delta
f(d)=-(\rho+2),\;\rho\ge-1 & \mathsf{upper} & \mathsf{Lem:}\ref%
{lemi:delta2-formula}(8) \\ \hline
19 & 2g(d-1)\le(\rho+2)(\rho+3),\quad\Delta g(d)\le0 & \mathsf{bound} & 
\mathsf{Eq}(7) \\ \hline
\end{array}%
$

\nointerlineskip\noindent$%
\begin{array}{|m{0.5cm}|l|r|c|}
\hline
20 & H(d-1)\le(\rho+2)(\rho+3)\wedge\Delta H(d)=-(\rho+2)+\Delta g(d) & 
\mathsf{combine} & \mathsf{Eq}(18,19) \\ \hline
21 & (\rho+2)(\rho+3)+2g(d-1)\le2\bigl(-(\rho+2)+\Delta g(d)\bigr)^{2} & 
\mathsf{algebra} & \mathsf{Eq}(19),\rho\ge-1 \\ \hline
22 & H(d-1)\le(\Delta H(d))^{2} & \mathsf{conclude} & \mathsf{Eq}(20,21) \\ 
\hline
\end{array}%
$

\noindent\textit{Sub-cases B1 and B2 ($\Delta^{2}g=1$).} Use the shift
reduction: if $\Delta^{2}H(d+1)=1$ and $H(d+1)\le(\Delta H(d+1))^{2}$ then 
\begin{equation*}
H(d-1)=H(d+1)-2\Delta H(d+1)+1\le(\Delta H(d+1))^{2}-2\Delta
H(d+1)+1=(\Delta H(d))^{2}. 
\end{equation*}
Cases B1 and B2 are the mirror of A1 and A2 with $f$ and $g$ exchanged.

\subsubsection*{Step 4: Case $\Delta^{2}H(d+1)=2$.}

\noindent$%
\begin{array}{|m{0.5cm}|l|r|c|}
\hline
23 & \Delta^{2}H(d+1)=2\Rightarrow\Delta^{2}f(d+1)=\Delta^{2}g(d+1)=1 & 
\mathsf{case} & \mathsf{Eq}(8) \\ \hline
24 & \text{L/U\ alternatives at}\;d+1:\;\text{(L,L)},\text{(L,U)},\text{(U,L)%
},\text{(U,U)} & \mathsf{dichotomy} & \mathsf{Lem:}\ref{lemi:delta2-dich}(7)
\\ \hline
\end{array}%
$

\nointerlineskip\noindent$%
\begin{array}{|m{0.5cm}|l|r|c|}
\hline
25 & \text{(L,L):\ violates}\;m_{2}(x)>m_{1}(x) & \mathsf{contradiction} & 
\mathsf{Eq}(3,24) \\ \hline
26 & \text{(L,U):\ violates\ shift\ ordering} & \mathsf{contradiction} & 
\mathsf{Eq}(3,24) \\ \hline
\end{array}%
$

\nointerlineskip\noindent$%
\begin{array}{|m{0.5cm}|l|r|c|}
\hline
27 & \text{(U,L):\ violates\ shift\ ordering} & \mathsf{contradiction} & 
\mathsf{Eq}(3,24) \\ \hline
28 & \text{(U,U):\;put}\;\rho:=e-d-1.\;\mathsf{Upper\ formulas:}%
\;2f(d)=(\rho+2)(\rho+3),\;\Delta f(d+1)=-(\rho+2) & \mathsf{upper} & 
\mathsf{Lem:}\ref{lemi:delta2-formula}(8) \\ \hline
\end{array}%
$

\nointerlineskip\noindent$%
\begin{array}{|m{0.5cm}|l|r|c|}
\hline
29 & 2g(d)=(\rho+2)(\rho+3),\quad\Delta g(d+1)=-(\rho+2) & \mathsf{upper} & 
\mathsf{Lem:}\ref{lemi:delta2-formula}(8) \\ \hline
30 & \text{From U at}\;d+1:\;\rho=e-d-1\ge-2,\;\text{so}\;\Delta
H(d+1)=-2(\rho+2)\le0 & \mathsf{bound} & \mathsf{Eq}(28,29) \\ \hline
\end{array}%
$

\nointerlineskip\noindent$%
\begin{array}{|m{0.5cm}|l|r|c|}
\hline
31 & 2H(d)=2(\rho+2)(\rho+3) & \mathsf{combine} & \mathsf{Eq}(28,29) \\ 
\hline
32 & (\Delta H(d+1))^{2}-2H(d)=2(\rho+2)(\rho+1)\ge0\;\text{(since}%
\;\rho\ge-2) & \mathsf{algebra} & \mathsf{Eq}(30,31) \\ \hline
\end{array}%
$

\nointerlineskip\noindent$%
\begin{array}{|m{0.5cm}|l|r|c|}
\hline
33 & \Delta^{2}H(d+1)=2\Rightarrow2H(d)\le(\Delta H(d+1))^{2} & \mathsf{%
conclude} & \mathsf{Eq}(32) \\ \hline
\end{array}%
$

\subsubsection*{Step 5: Conclusion.}

\noindent From `$\Delta^{2}H(d+1)=2$' and `$H(d+1)\le(\Delta H(d+1))^{2}$'
(obtained in Step 4), apply the shift reduction: \noindent$%
\begin{array}{|m{0.5cm}|l|r|c|}
\hline
34 & H(d-1)=H(d+1)-2\Delta H(d+1)+2\;\mathsf{when}\;\Delta^{2}H(d+1)=2 & 
\mathsf{shift} & \mathsf{Eq}(30) \\ \hline
35 & \Delta H(d+1)=-2(\rho+2)\le0\;\text{(from}\;\rho\ge-2\text{)},\;\text{so%
}\;\Delta H(d+1)\le1 & \mathsf{bound} & \mathsf{Eq}(30) \\ \hline
\end{array}%
$

\nointerlineskip\noindent$%
\begin{array}{|m{0.5cm}|l|r|c|}
\hline
36 & H(d-1)\le(\Delta H(d+1))^{2}-2\Delta H(d+1)+2\le(\Delta
H(d+1)-2)^{2}=(\Delta H(d))^{2} & \mathsf{algebra} & \mathsf{Eq}(32,34,35)
\\ \hline
37 & H(d-1)H(d+1)\le H(d)^{2}\;\mathsf{for\;all}\;0<d<e & \mathsf{combine} & 
\mathsf{Eq}(9,10,11,17,22,36;\text{B1/B2}) \\ \hline
38 & \mathsf{LC}(\mathsf{socle}(P),h_{P}) & \mathsf{conclude} & \mathsf{Eq}%
(37) \\ \hline
39 & \top & \mathsf{close} & \to\mathsf{I}(1,\dots,38) \\ \hline
\end{array}%
$

\medskip

% ------------------------------------------------------------

\subsection*{Theorem \protect\ref{thmi:P6} (Open)}

% ------------------------------------------------------------
\noindent\textbf{Status: open.} Theorem~P6 (unimodality of type-$2$ pure $O$%
-sequences in codimension $r\ge5$) is stated in the Statements section for
reference but is \emph{not} proved in this document. The proofs of P4 and P5
rely essentially on the codimension-$3$ structure (three coordinates,
allowing a single dominant coordinate $x$ to be isolated and the swap
reduction to be applied). This structure does not directly generalise to $%
r\ge5$, and the corresponding question was identified as open by Zanello~%
\cite{zanello2024}. No proof table is provided for P6.

\medskip

% ------------------------------------------------------------

\subsection*{Proof of Theorem \protect\ref{thmi:P7} (Eventual
Non-Log-Concavity)}

% ------------------------------------------------------------
\noindent$%
\begin{array}{|m{0.5cm}|l|r|c|}
\hline
1 & r\ge3,\;t\ge5 & \mathsf{hypothesis} & \mathsf{Hyp} \\ \hline
2 & m_{\mathrm{fat}}=x^{t}y^{t}z^{t},\;m_{i}=x^{3t-i}y^{i}\;(0\le i<t-1) & 
\mathsf{construction} & \mathsf{Def:}\ref{defi:Monomial}(1) \\ \hline
\end{array}%
$

\nointerlineskip\noindent$%
\begin{array}{|m{0.5cm}|l|r|c|}
\hline
3 & \Gamma_{t}:=\Gamma(m_{\mathrm{fat}},m_{0},\dots,m_{t-2}) & \mathsf{ideal}
& \mathsf{Def:}\ref{defi:generatedOrderIdeal}(4) \\ \hline
4 & \mathsf{pure},\;\mathsf{socle}=3t,\;\mathsf{type}=t & \mathsf{from\ Lemma%
} & \mathsf{Lem:}\ref{lemi:exists-non-lc-r3}(5,6,7) \\ \hline
\end{array}%
$

\nointerlineskip\noindent$%
\begin{array}{|m{0.5cm}|l|r|c|}
\hline
5 & h_{\Gamma_{t}}(3t)=t,\;h_{\Gamma_{t}}(3t-1)=t+2,\;h_{%
\Gamma_{t}}(3t-2)=t+5 & \mathsf{count\ (\S above)} & \mathsf{Lem:}\ref%
{lemi:exists-non-lc-r3}(8,9,10) \\ \hline
6 & (t+5)\cdot t>(t+2)^{2}\;\mathsf{for}\;t\ge5 & \mathsf{algebra} & \mathsf{%
Eq}(5) \\ \hline
\end{array}%
$

\nointerlineskip\noindent$%
\begin{array}{|m{0.5cm}|l|r|c|}
\hline
7 & \neg\mathsf{LC}(3t,h_{\Gamma_{t}}) & \mathsf{conclude} & \mathsf{Lem:}%
\ref{lemi:exists-non-lc-r3}(12) \\ \hline
8 & \exists P:\mathsf{PureOrderIdeal}_{r},\;\mathsf{type}(P)=t\wedge\neg%
\mathsf{LC}(P) & \mathsf{embed} & \mathsf{Lem:}\ref{lemi:embed-preserves}(7)
\\ \hline
\end{array}%
$

\nointerlineskip\noindent$%
\begin{array}{|m{0.5cm}|l|r|c|}
\hline
9 & \exists T\ge1,\;\forall t\ge T,\;\mathsf{existsNonLC}(r,t) & \mathsf{take%
}\;T:=5 & \mathsf{Eq}(8) \\ \hline
10 & \top & \mathsf{close} & \to\mathsf{I}(1,\dots,9) \\ \hline
\end{array}%
$

\medskip

% ------------------------------------------------------------

\subsection*{Proof of Theorem \protect\ref{thmi:P8} (Least Threshold)}

% ------------------------------------------------------------
\noindent$%
\begin{array}{|m{0.5cm}|l|r|c|}
\hline
1 & r\ge3 & \mathsf{hypothesis} & \mathsf{Hyp} \\ \hline
2 & \exists T\ge1,\;\forall t\ge T,\;\mathsf{existsNonLC}(r,t) & \mathsf{P7}
& \mathsf{Thm:}\ref{thmi:P7}(9) \\ \hline
\end{array}%
$

\nointerlineskip\noindent$%
\begin{array}{|m{0.5cm}|l|r|c|}
\hline
3 & S:=\{t\in\N\mid\forall t^{\prime}\ge t,\;\mathsf{existsNonLC}%
(r,t^{\prime})\} & \mathsf{set} & \mathsf{Eq}(2) \\ \hline
4 & S\ne\emptyset & \mathsf{from}\;T & \mathsf{Eq}(2,3) \\ \hline
\end{array}%
$

\nointerlineskip\noindent$%
\begin{array}{|m{0.5cm}|l|r|c|}
\hline
5 & t_{0}:=\mathsf{IsLeast\text{-}witness}(S) & \mathsf{least-threshold} & 
\mathsf{Lem:}\ref{lemi:least-threshold}(6) \\ \hline
6 & \mathsf{IsLeast}(t_{0},S) & \mathsf{conclude} & \mathsf{Eq}(4,5) \\ 
\hline
\end{array}%
$

\nointerlineskip\noindent$%
\begin{array}{|m{0.5cm}|l|r|c|}
\hline
7 & \top & \mathsf{close} & \to\mathsf{I}(1,\dots,6) \\ \hline
\end{array}%
$

\section{Scripts}

\label{sec:scripts} 
% ============================================================
\noindent Each tree records, for one result, the entries used by its proof,
transitively closed. Names coincide with the names that appear in the
statement tables. Leaves are either postulates (unproved classical inputs),
axioms of \textsf{ZF+DC}, or primitive definitions. In the trees below the
prefix \verb|Ld | carries the depth of the node (root at \verb|L0|) and the
name is indented \verb|2*d| spaces after the pipe, so every chain can be
read straight down a column of increasing depth labels.

\medskip \noindent Every tree has a corresponding Lean~4 formalization in
one of the files \texttt{hilbert\_functions\_1.lean} through \texttt{%
hilbert\_functions\_8.lean}; the table below records the mapping from the
entries of the Statements section to the Lean file and to the top-level Lean
declaration that discharges the proof obligation.%
\begin{equation*}
\begin{array}{|m{3.5cm}|m{5cm}|m{5.5cm}|}
\hline
\text{\textsf{Entry}} & \text{\textbf{Lean file}} & \text{\textsf{Lean
declaration}} \\ \hline
\text{\textsf{principalOSeq}} & \text{\texttt{hilbert\_functions\_2.lean}} & 
\text{\texttt{unimodal\_seq\_of\_g}} \\ \hline
\text{\textsf{principalPos}} & \text{\texttt{hilbert\_functions\_2.lean}} & 
\text{\texttt{seq\_of\_g\_pos}} \\ \hline
\text{\textsf{twoMaxDecomp}} & \text{\texttt{hilbert\_functions\_4.lean}} & 
\text{\texttt{carrier\_of\_maximals}} \\ \hline
\text{\textsf{shiftReduce}} & \text{\texttt{hilbert\_functions\_4.lean}} & 
\text{\texttt{diff\_eq\_shifted\_OSeq}} \\ \hline
\text{\textsf{delta2Dich}} & \text{\texttt{hilbert\_functions\_5.lean}} & 
\text{\texttt{diff2\_h\_seq\_eq\_one\_implies}} \\ \hline
\text{\textsf{LUDich}} & \text{\texttt{hilbert\_functions\_5.lean}} & \text{%
\texttt{diff2\_H\_seq\_le\_two}} \\ \hline
\text{\textsf{existsNonLC\_r3}} & \text{\texttt{hilbert\_functions\_7.lean}}
& \text{\texttt{dip\_exists\_dip}} \\ \hline
\text{\textsf{embedPres}} & \text{\texttt{hilbert\_functions\_7.lean}} & 
\text{\texttt{reduction\_to\_r3}} \\ \hline
\text{\textsf{leastThreshold}} & \text{\texttt{hilbert\_functions\_8.lean}}
& \text{\texttt{exists\_isLeast\_of\_exists}} \\ \hline
\text{\textsf{P1}} & \text{\texttt{hilbert\_functions\_1.lean}} & \text{%
\texttt{P1}} \\ \hline
\text{\textsf{P2}} & \text{\texttt{hilbert\_functions\_2.lean}} & \text{%
\texttt{P2}} \\ \hline
\text{\textsf{P3}} & \text{\texttt{hilbert\_functions\_3.lean}} & \text{%
\texttt{P3}} \\ \hline
\text{\textsf{P4}} & \text{\texttt{hilbert\_functions\_4.lean}} & \text{%
\texttt{P4}} \\ \hline
\text{\textsf{P5}} & \text{\texttt{hilbert\_functions\_5.lean}} & \text{%
\texttt{P5}} \\ \hline
\text{\textsf{P7}} & \text{\texttt{hilbert\_functions\_7.lean}} & \text{%
\texttt{P7}} \\ \hline
\text{\textsf{P8}} & \text{\texttt{hilbert\_functions\_8.lean}} & \text{%
\texttt{P8}} \\ \hline
\end{array}%
\end{equation*}

\medskip

\subsection*{Dependency tree: \textsf{LCcore}}

\noindent%
\begin{minipage}[t]{0.48\linewidth}
	\vspace{0pt}
	\begin{verbatimbox}
		L0 | LC-core-alg
		L1 |   Def:Z (integers)
	\end{verbatimbox}
\end{minipage}\hfill 
\begin{minipage}[t]{0.48\linewidth}
	\vspace{0pt}
	\textbf{Annotation.} Lean proof: \texttt{log\_concave\_core}, by trichotomy on $z$ and elementary arithmetic.
\end{minipage}

\medskip

\subsection*{Dependency tree: \textsf{delta2Formula}}

\noindent%
\begin{minipage}[t]{0.48\linewidth}
	\vspace{0pt}
	\begin{verbatimbox}
		L0 | delta2-formula
		L1 |   Def:divisorsFinset
		L1 |   Def:totalDeg
		L1 |   Def:Monomial
	\end{verbatimbox}
\end{minipage}\hfill 
\begin{minipage}[t]{0.48\linewidth}
	\vspace{0pt}
	\textbf{Annotation.} Lean proof: \texttt{diff2\_h\_seq\_eq\_Z\_calc}, by inclusion-exclusion and second-difference computation.
\end{minipage}

\medskip

\subsection*{Dependency tree: \textsf{principalOSeq}}

\noindent%
\begin{minipage}[t]{0.48\linewidth}
	\vspace{0pt}
	\begin{verbatimbox}
		L0 | principal-OSeq
		L1 |   pos:conv-LC
		L1 |   Def:IsLogConcaveUpTo
		L1 |   Def:totalDeg
		L1 |   Def:divisorsFinset
	\end{verbatimbox}
\end{minipage}\hfill 
\begin{minipage}[t]{0.48\linewidth}
	\vspace{0pt}
	\textbf{Annotation.} Lean proof: \texttt{unimodal\_seq\_of\_g}, via \texttt{Finsupp.induction}. The convolution step uses the ratio-decreasing characterization of log-concavity proved locally as \texttt{log\_concave\_conv\_Z} and \texttt{is\_ratio\_decreasing\_Z}.
\end{minipage}

\medskip

\subsection*{Dependency tree: \textsf{principalPos}}

\noindent%
\begin{minipage}[t]{0.48\linewidth}
	\vspace{0pt}
	\begin{verbatimbox}
		L0 | principal-pos
		L1 |   Def:Monomial
		L1 |   Def:totalDeg
		L1 |   Def:divisorsFinset
	\end{verbatimbox}
\end{minipage}\hfill 
\begin{minipage}[t]{0.48\linewidth}
	\vspace{0pt}
	\textbf{Annotation.} Lean proof: \texttt{seq\_of\_g\_pos} via the truncation map $m_{i}\mapsto\min(g_{i},d-\sum_{j<i}m_{j})$.
\end{minipage}

\medskip

\subsection*{Dependency tree: \textsf{twoMaxDecomp}}

\noindent%
\begin{minipage}[t]{0.48\linewidth}
	\vspace{0pt}
	\begin{verbatimbox}
		L0 | two-max-decomp
		L1 |   Def:typePO
		L1 |   Def:IsMaximalIn
		L1 |   Def:OrderIdeal
		L1 |   Def:divisorsFinset
	\end{verbatimbox}
\end{minipage}\hfill 
\begin{minipage}[t]{0.48\linewidth}
	\vspace{0pt}
	\textbf{Annotation.} Lean proof: \texttt{carrier\_of\_maximals}, which uses \texttt{Finset.exists\_max\_image} to terminate every chain in $\Gamma$ at one of the two maximal monomials.
\end{minipage}

\medskip

\subsection*{Dependency tree: \textsf{shiftReduce}}

\noindent%
\begin{minipage}[t]{0.48\linewidth}
	\vspace{0pt}
	\begin{verbatimbox}
		L0 | shift-reduce
		L1 |   two-max-decomp
		L2 |     Def:typePO
		L1 |   Def:divisorsFinset
		L1 |   Def:Monomial
	\end{verbatimbox}
\end{minipage}\hfill 
\begin{minipage}[t]{0.48\linewidth}
	\vspace{0pt}
	\textbf{Annotation.} Lean proof: \texttt{OSeq\_union} combined with \texttt{monomialInf} (the coordinatewise min) and \texttt{diff\_eq\_shifted\_OSeq}.
\end{minipage}

\medskip

\subsection*{Dependency tree: \textsf{delta2Dich}}

\noindent%
\begin{minipage}[t]{0.48\linewidth}
	\vspace{0pt}
	\begin{verbatimbox}
		L0 | delta2-dich
		L1 |   Lem:delta2-formula
		L2 |     Def:totalDeg
		L2 |     Def:Monomial
	\end{verbatimbox}
\end{minipage}\hfill 
\begin{minipage}[t]{0.48\linewidth}
	\vspace{0pt}
	\textbf{Annotation.} Lean proof: \texttt{diff2\_h\_seq\_eq\_one\_implies}. The L/U dichotomy is extracted from the twelve-term indicator sum by \texttt{grind} on \texttt{Fin 3}.
\end{minipage}

\medskip

\subsection*{Dependency tree: \textsf{LUDich}}

\noindent%
\begin{minipage}[t]{0.48\linewidth}
	\vspace{0pt}
	\begin{verbatimbox}
		L0 | LU-dich
		L1 |   delta2-dich
		L2 |     Lem:delta2-formula
		L1 |   shift-reduce
		L2 |     two-max-decomp
	\end{verbatimbox}
\end{minipage}\hfill 
\begin{minipage}[t]{0.48\linewidth}
	\vspace{0pt}
	\textbf{Annotation.} Lean proof: \texttt{diff2\_H\_seq\_le\_two}, which feeds \texttt{diff2\_H\_seq\_eq\_add} into \texttt{diff2\_h\_seq\_le\_one} twice.
\end{minipage}

\medskip

\subsection*{Dependency tree: \textsf{existsNonLC\_r3}}

\noindent%
\begin{minipage}[t]{0.48\linewidth}
	\vspace{0pt}
	\begin{verbatimbox}
		L0 | exists-non-lc-r3
		L1 |   Def:Monomial
		L1 |   Def:generatedOrderIdeal
		L1 |   Def:pureOSequence
		L1 |   Def:IsLogConcaveUpTo
	\end{verbatimbox}
\end{minipage}\hfill 
\begin{minipage}[t]{0.48\linewidth}
	\vspace{0pt}
	\textbf{Annotation.} Lean proof: \texttt{dip\_exists\_dip} together with \texttt{exists\_dip\_sequence\_r3}. The construction is \texttt{dipPureOrderIdeal3}.
\end{minipage}

\medskip

\subsection*{Dependency tree: \textsf{embedPres}}

\noindent%
\begin{minipage}[t]{0.48\linewidth}
	\vspace{0pt}
	\begin{verbatimbox}
		L0 | embed-pres
		L1 |   exists-non-lc-r3
		L2 |     Def:Monomial
		L2 |     Def:generatedOrderIdeal
		L1 |   Def:totalDeg
	\end{verbatimbox}
\end{minipage}\hfill 
\begin{minipage}[t]{0.48\linewidth}
	\vspace{0pt}
	\textbf{Annotation.} Lean proof: \texttt{reduction\_to\_r3}. The embedding $\iota:\N^{3}\hookrightarrow\N^{r}$ is \texttt{embedMonomial} / \texttt{embedPureOrderIdeal}.
\end{minipage}

\medskip

\subsection*{Dependency tree: \textsf{leastThreshold}}

\noindent%
\begin{minipage}[t]{0.48\linewidth}
	\vspace{0pt}
	\begin{verbatimbox}
		L0 | least-threshold
		L1 |   ax:DC
		L1 |   Nat.wellFounded
	\end{verbatimbox}
\end{minipage}\hfill 
\begin{minipage}[t]{0.48\linewidth}
	\vspace{0pt}
	\textbf{Annotation.} Lean proof: \texttt{exists\_isLeast\_of\_exists}, which reduces to \texttt{Nat.find} and \texttt{Nat.find\_min'}.
\end{minipage}

\medskip

\subsection*{Dependency tree: \textsf{P1}}

\noindent%
\begin{minipage}[t]{0.48\linewidth}
	\vspace{0pt}
	\begin{verbatimbox}
		L0 | P1
		L1 |   Def:generatedOrderIdeal
		L1 |   Def:pureOSequence
	\end{verbatimbox}
\end{minipage}\hfill 
\begin{minipage}[t]{0.48\linewidth}
	\vspace{0pt}
	\textbf{Annotation.} Lean proof: \texttt{P1}, splitting into \texttt{GammaP1\_0}, \texttt{GammaP1\_1}, \texttt{GammaP1\_2}, \texttt{GammaP1\_3}.
\end{minipage}

\medskip

\subsection*{Dependency tree: \textsf{P2}}

\noindent%
\begin{minipage}[t]{0.48\linewidth}
	\vspace{0pt}
	\begin{verbatimbox}
		L0 | P2
		L1 |   principal-OSeq
		L2 |     pos:conv-LC
		L1 |   principal-pos
		L1 |   Def:typePO
		L1 |   Def:IsLogConcaveUpTo
		L1 |   Def:IsUnimodal
	\end{verbatimbox}
\end{minipage}\hfill 
\begin{minipage}[t]{0.48\linewidth}
	\vspace{0pt}
	\textbf{Annotation.} Lean proof: \texttt{P2}, composed of \texttt{P2\_coeff}, \texttt{P2\_unimodal}, and \texttt{P2\_logconcave}.
\end{minipage}

\medskip

\subsection*{Dependency tree: \textsf{P3}}

\noindent%
\begin{minipage}[t]{0.48\linewidth}
	\vspace{0pt}
	\begin{verbatimbox}
		L0 | P3
		L1 |   pos:Macaulay2
		L1 |   Def:pureOSequence
		L1 |   Def:OrderIdeal
		L1 |   Def:IsLogConcaveUpTo
	\end{verbatimbox}
\end{minipage}\hfill 
\begin{minipage}[t]{0.48\linewidth}
	\vspace{0pt}
	\textbf{Annotation.} Lean proof: \texttt{P3}, via \texttt{P3\_dichotomy} on the slice sizes \texttt{slice\_size} and \texttt{macaulay\_2\_vars}. The dichotomy is closed by \texttt{nat\_AMGM}.
\end{minipage}

\medskip

\subsection*{Dependency tree: \textsf{P4}}

\noindent%
\begin{minipage}[t]{0.48\linewidth}
	\vspace{0pt}
	\begin{verbatimbox}
		L0 | P4
		L1 |   two-max-decomp
		L1 |   shift-reduce
		L2 |     two-max-decomp
		L1 |   principal-OSeq
		L2 |     pos:conv-LC
		L1 |   Def:IsUnimodal
	\end{verbatimbox}
\end{minipage}\hfill 
\begin{minipage}[t]{0.48\linewidth}
	\vspace{0pt}
	\textbf{Annotation.} Lean proof: \texttt{P4}, using \texttt{unimodal\_of\_two\_maximals} which packages \texttt{oseq\_union\_eq\_add\_diff} and \texttt{unimodal\_of\_OSeq\_add\_shifted}.
\end{minipage}

\medskip

\subsection*{Dependency tree: \textsf{P5}}

\noindent%
\begin{minipage}[t]{0.48\linewidth}
	\vspace{0pt}
	\begin{verbatimbox}
		L0 | P5
		L1 |   two-max-decomp
		L1 |   shift-reduce
		L2 |     two-max-decomp
		L1 |   delta2-dich
		L2 |     Lem:delta2-formula
		L1 |   LU-dich
		L2 |     delta2-dich
		L1 |   Lem:LC-core-alg
		L1 |   Def:IsLogConcaveUpTo
	\end{verbatimbox}
\end{minipage}\hfill 
\begin{minipage}[t]{0.48\linewidth}
	\vspace{0pt}
	\textbf{Annotation.} Lean proof: \texttt{P5}, whose heart is \texttt{sum\_log\_concave\_helper}, composed of \texttt{log\_concave\_core}, \texttt{diff2\_H\_seq\_le\_two}, and the four case-lemmas \texttt{diff2\_H\_seq\_eq\_one\_implies\_inl} / \texttt{\_inr} and \texttt{diff2\_H\_seq\_eq\_two\_implies}.
\end{minipage}

\medskip

\subsection*{Dependency tree: \textsf{P7}}

\noindent%
\begin{minipage}[t]{0.48\linewidth}
	\vspace{0pt}
	\begin{verbatimbox}
		L0 | P7
		L1 |   exists-non-lc-r3
		L2 |     Def:Monomial
		L2 |     Def:generatedOrderIdeal
		L2 |     Def:pureOSequence
		L1 |   embed-pres
		L2 |     exists-non-lc-r3
		L2 |     Def:totalDeg
		L1 |   Def:ExistsNonLogConcave
	\end{verbatimbox}
\end{minipage}\hfill 
\begin{minipage}[t]{0.48\linewidth}
	\vspace{0pt}
	\textbf{Annotation.} Lean proof: \texttt{P7}, taking $T:=5$. The construction is \texttt{dipPureOrderIdeal3}.
\end{minipage}

\medskip

\subsection*{Dependency tree: \textsf{P8}}

\noindent%
\begin{minipage}[t]{0.48\linewidth}
	\vspace{0pt}
	\begin{verbatimbox}
		L0 | P8
		L1 |   P7
		L2 |     exists-non-lc-r3
		L2 |     embed-pres
		L1 |   least-threshold
		L2 |     ax:DC
		L1 |   Def:ExistsNonLogConcave
	\end{verbatimbox}
\end{minipage}\hfill 
\begin{minipage}[t]{0.48\linewidth}
	\vspace{0pt}
	\textbf{Annotation.} Lean proof: \texttt{P8}, which is \texttt{P7\_implies\_P8} applied to the threshold existence of P7.
\end{minipage}

\medskip

\noindent\textbf{Theorem P6 (open).} No tree is given for P6; no Lean file
among the seven supplied formalizes a proof; P6 is recorded as open in the
Statements section.

\medskip

\subsubsection*{Complete leaves}

\noindent%
\begin{minipage}[t]{0.48\linewidth}
	\vspace{0pt}
	\textbf{Postulates.}
	\smallskip
	\begin{verbatimbox}
		pos:conv-LC
		pos:Macaulay2
	\end{verbatimbox}
\end{minipage}\hfill 
\begin{minipage}[t]{0.48\linewidth}
	\vspace{0pt}
	\textbf{Axioms.}
	\smallskip
	\begin{verbatimbox}
		ax:DC
	\end{verbatimbox}
\end{minipage}

\medskip

\noindent%
\begin{minipage}[t]{\linewidth}
	\vspace{0pt}
	\textbf{Definitions used as leaves.}
	\smallskip
	\begin{verbatimbox}
		Def:Monomial        Def:totalDeg       Def:OrderIdeal
		Def:divisorsFinset  Def:generatedOrderIdeal
		Def:IsMaximalIn     Def:maximalMonomials
		Def:IsPure          Def:PureOrderIdeal Def:typePO
		Def:socleDegree     Def:pureOSequence  Def:IsUnimodal
		Def:IsLogConcaveUpTo Def:ExistsNonLogConcave
	\end{verbatimbox}
\end{minipage}

% ============================================================

\section{Verification}

\label{sec:verification} 
% ============================================================
\noindent\textbf{Formal verification.} Every lemma and theorem in the
Statements section (with the exception of the open Theorem~P6) has a
corresponding machine-checked proof in the Lean~4 proof assistant, in one of
the files \texttt{hilbert\_functions\_1.lean} through \texttt{%
hilbert\_functions\_8.lean}. The Scripts section above records the mapping
from each entry to its Lean file and to the top-level Lean declaration that
discharges the proof obligation. All Lean proofs are publicly available at %
\url{https://github.com/google-deepmind/alphaproof-nexus-results}.

\medskip\noindent\textbf{Verification methodology.} The Lean formalization
proceeds by:

\begin{enumerate}
\item \textbf{Formalization of the statement.} Each theorem is stated in
Lean~4 as a \texttt{Prop} involving the definitions of the Definitions
section, following the conventions used in the Mathlib library \cite%
{mathlib2020}.

\item \textbf{Proof search.} The proofs were found by the AlphaProof Nexus
agent \cite{tsoukalas2026}, which searches for tactic scripts in the Lean~4
tactic language, guided by compiler feedback.

\item \textbf{Machine verification.} The Lean compiler checks every proof
step; a theorem is accepted only if the compiler reaches a state with no
remaining goals.

\item \textbf{Informal reading.} The step tables in the Proofs section are
an \emph{informal} deformation of the Lean proofs: they track the same
logical structure but compress many Lean-level tactic steps into single
mathematical assertions. The Lean files are the authoritative formal
artefact; the tables here are a reading aid.
\end{enumerate}

\medskip\noindent\textbf{Robustness against known failure modes.} The
AlphaProof Nexus preprint \cite{tsoukalas2026} documents two failure modes
of LLM-based proof search that our verification workflow explicitly guards
against:

\begin{itemize}
\item \emph{Sorried core lemmas.} The Lean proof files in the repository
contain no \texttt{sorry} in any of the declarations listed in the Scripts
mapping table. A `sorry` audit at the declaration level---reading off the
declaration names and confirming each is a closed proof term---is the basic
check that a reader can run without any Lean expertise.

\item \emph{Hallucinated literature citations.} The preprint reports that in
its failure cases, top sketches relied on lemmas marked \texttt{sorry} that
the agent claimed were established results. Since all sixteen mapping-table
declarations are themselves proved in Lean (not deferred to cited work),
this failure mode does not affect the results presented here.
\end{itemize}

\medskip\noindent\textbf{Known limitations.} The Lean formalization verifies
the formal statements as given; the informal reading of those statements
against the original conjectures of Stanley,
Boij--Migliore--Mir\'o-Roig--Nagel--Zanello, and Zanello is done by hand and
is not machine-checked. The deformation tables in the Proofs section are not
themselves verified; they should be read alongside the Lean source.

\medskip\noindent\textbf{Reproducibility.} The Lean development builds
against Mathlib \cite{mathlib2020}. To reproduce the machine check of any
specific declaration (e.g.\ \texttt{P5}), clone the repository, check out
the pinned Mathlib commit, and run \texttt{lake build} on the relevant 
\texttt{hilbert\_functions\_*.lean} file. The declaration \texttt{P5}
compiles successfully and its type is \texttt{ProblemP5}, which unfolds to
the statement of P5 in the Statements section.

\medskip\noindent\textbf{Status.} The results described here are reported in 
\cite{tsoukalas2026}, a preprint of Google DeepMind's AlphaProof Nexus
project. The work has not yet appeared in a peer-reviewed venue at the time
of writing.

% ============================================================

\section{Discussion and Limitations}

% ============================================================
\noindent\textbf{Dependence on formal verification.} The mathematical weight
of this document rests on the Lean formalization. The Proofs section is a 
\emph{reading aid}, not a self-contained proof: the deepest sub-arguments
(the L/U case analysis behind \texttt{delta2Dich}, the first-half
log-concavity step in P3, the four-case analysis in P5) are presented at a
level of detail consistent with a mathematical exposition, not a formal
proof. A reader who wishes to independently verify the results should
consult the Lean files.

\medskip\noindent\textbf{Two postulates.} Only two postulates are retained: 
\texttt{conv-LC} (the convolution theorem for log-concave sequences) and 
\texttt{Macaulay2} (Macaulay's two-variable theorem). Both are classical and
well-documented in the literature (see \S \ref{sec:verification} and the
Provenance note in \S 3). They are admitted here to keep the exposition
focused on the eight-problem layer.

\medskip\noindent\textbf{The open question.} P6 (unimodality of type-$2$
pure $O$-sequences in codimension $r \ge 5$) is stated but not resolved. The
proofs of P4 and P5 depend essentially on the codimension-$3$ structure, in
particular the existence of a single dominant coordinate $x$ that allows the
swap reduction. This does not generalise to $r \ge 5$, and the corresponding
question remains open.

\medskip\noindent\textbf{Role of AI in the proofs.} The Lean proofs were
found by an LLM-driven agent (\cite{tsoukalas2026}), not written by hand.
This is a new mode of mathematical production, and it has two consequences
the reader should be aware of:

\begin{itemize}
\item The \emph{correctness} of each proof is assured by the Lean compiler;
the question of whether the formalized statement matches the informal
conjecture is a separate, human-audited step (see Verification above).

\item The \emph{mathematical insight} behind each proof is not necessarily
reflected in the agent's internal reasoning; the natural language summaries
in this document are reverse-engineered from the Lean proofs.
\end{itemize}

% ============================================================

\section{Conclusion}

% ============================================================
\noindent The document presents the dependency-ordered exposition of eight
problems on pure $O$-sequences, organised into six parts (Axioms,
Postulates, Definitions, Statements, Proofs, Scripts), with the choice
strength of every result recorded as \textsf{ZF}, \textsf{ZF+DC}, or \textsf{%
ZFC}. Within the Statements part, Lemmas and Theorems form a single
dependency layer, so long as each entry appears only after all entries it
depends on. In particular \textsf{LCcore} and \textsf{delta2Formula} are
stated first among the Lemmas because \textsf{delta2Dich} and \textsf{P5}
depend on them; \textsf{principal-OSeq} is stated because P2 depends on it; 
\textsf{two-max-decomp} and \textsf{shift-reduce} precede P4 and P5; \textsf{%
delta2-dich} and \textsf{LU-dich} precede P5; and \textsf{exists-non-lc-r3}
precedes P7 and P8.

Only two postulates remain (\textsf{conv-LC} and \textsf{Macaulay2}), both
cited to the literature. The codimension-$r \ge 5$ type-$2$ unimodality
question (P6) is stated but not resolved. Every lemma and theorem stated in
the Statements section has a corresponding proof table in the Proofs section
(except P6), and the proofs are presented in the same order as the
statements. The Proofs section is deliberately an \emph{informal deformation}
of the machine-checked Lean proofs; readers who wish to verify the results
mechanically should consult the Lean files in the companion repository and
read the Verification section for methodology, robustness checks, and
limitations.

% ============================================================

\begin{thebibliography}{99}
\bibitem{stanley1978} R.~P. Stanley. \newblock Hilbert functions of graded 
algebras. \newblock \emph{Adv.\ Math.} \textbf{28} (1978), 57--83. 
\newblock
\url{https://doi.org/10.1016/0001-8708(78)90045-2}. \medskip

\bibitem{bmmnz2012} M.~Boij, J.~C. Migliore, R.~M. Mir\'o-Roig, U.~Nagel, 
F.~Zanello. \newblock \emph{On the Shape of a Pure $O$-Sequence}. \newblock
Mem.\ Amer.\ Math.\ Soc.\ \textbf{218} (2012), no.~1024. \newblock
\url{https://doi.org/10.1090/S0065-9266-2011-00647-7}. \medskip

\bibitem{zanello2024} F.~Zanello. \newblock Log-concavity of level Hilbert 
functions and pure $O$-sequences. \newblock \emph{J.\ Commut.\ Algebra}  
\textbf{16} (2024), 245--256. \newblock
\url{https://doi.org/10.1216/jca.2024.16.245}. \medskip

\bibitem{tsoukalas2026} G.~Tsoukalas, A.~Kovsharov, S.~Shirobokov, et~al.  %
\newblock Advancing mathematics research with AI-driven formal proof search.
\newblock arXiv:2605.22763 (2026). \newblock
\url{https://arxiv.org/abs/2605.22763}. \medskip

\bibitem{simpson2009} S.~G. Simpson. \newblock \emph{Subsystems of Second
Order Arithmetic}, 2nd ed. \newblock Perspectives in Logic, Cambridge 
University Press, Cambridge, 2009. \newblock
\url{https://doi.org/10.1017/CBO9780511581007}. \medskip

\bibitem{karlin1968} S.~Karlin. \newblock \emph{Total Positivity}, Vol.~I.  %
\newblock Stanford University Press, Stanford, CA, 1968. \newblock
\url{https://archive.org/details/totalpositivity0000karl}. \medskip

\bibitem{liggett1997} T.~M. Liggett. \newblock Ultra log-concavity and 
down-up chains. \newblock \emph{J.\ Math.\ Anal.\ Appl.} \textbf{213} 
(1997), 561--574. \newblock
\url{https://doi.org/10.1006/jmaa.1997.5581}. \medskip

\bibitem{macaulay1927} F.~S. Macaulay. \newblock Some properties of 
enumeration functions in the theory of modular systems. \newblock
\emph{J.\ London Math.\ Soc.} \textbf{2} (1927), 164--176. \newblock
\url{https://doi.org/10.1112/jlms/s1-2.3.164}. \medskip

\bibitem{brunsherzog1993} W.~Bruns, H.~J. Herzog. \newblock \emph{Cohen--
Macaulay Rings}. \newblock Cambridge Studies in Advanced Mathematics, 
vol.~39, Cambridge University Press, Cambridge, 1993. \newblock
\url{https://doi.org/10.1017/CBO9781139172073}. \medskip

\bibitem{mathlib2020} {The Mathlib Community}. \newblock The Lean 
mathematical library. \newblock In \emph{Proceedings of the 9th ACM SIGPLAN
International Conference on Certified Programs and Proofs},  CPP~2020, ACM,
New York, 2020. \newblock
\url{https://doi.org/10.1145/3372885.3373824}.
\end{thebibliography}

\end{document}
